% Part II: The Judgment Framework
% This file is \include'd from main.tex

\chapter{Design Principles --- Why Four Objects}
\label{ch:motivation}

Every judgment answers a single question:
\begin{quote}
  \emph{Over this space, how well can we construct a global solution,
  and what prevents it?}
\end{quote}
This question decomposes into exactly four irreducible aspects:
\begin{itemize}[nosep]
\item \textbf{Where} --- the context, the space over which the
  judgment is made;
\item \textbf{What} --- the construction, the content being asserted;
\item \textbf{Why not} --- the obstructions that prevent local pieces
  from assembling into a global whole;
\item \textbf{How good} --- the quality of the construction, measured
  in a graded algebraic structure.
\end{itemize}
Each aspect corresponds to a deep, well-studied branch of mathematics:
topos theory, type theory, cohomology, and quantale theory.  This
chapter explains why a judgment naturally has exactly four aspects
and why each demands its own mathematical universe.

\section{Why structured objects, not flat components}

When judgment components are listed as independent entries in a flat
tuple, their intricate dependencies must be imposed as external
axioms.  But the components of a judgment are not independent: the
coordinate~$c$ determines which
propositions~$\varphi$ make sense; the type~$A$ constrains which
terms~$E$ can serve as evidence; the cover encoded in~$O$ determines
which cohomology classes~$B$ can arise; the grade~$T$ summarises the
trail~$\Pi$.

Encoding these dependencies externally is
possible but unnatural.  It is analogous to defining a group as a
set~$G$ with a binary operation~$\mu$, an identity element~$e$, and an
inversion map~$\iota$, then imposing associativity, identity, and
inverse axioms externally.  The definition \emph{works}, but it hides
the fact that the axioms are not arbitrary constraints---they are the
defining properties of a single algebraic concept.

The 4-tuple formulation avoids this: each of the four objects
\emph{internalises} the relevant constraints, just as a group
internalises its axioms.

\section{The first aspect: $(c,\varphi)$ is a presheaf evaluation}
\label{sec:obs-topos}

Let $C$ be the site category whose objects are the coordinates of the
judgment, equipped with a Grothendieck topology~$J$ encoding the
coverage relation.  A proposition~$\varphi$ defined at coordinate~$c$
is a section of a presheaf $\varphi \in \ob(\PSh(C))$ evaluated at the
object~$c$:
\[
  \varphi(c) \;\in\; \Set.
\]
The pair $(c,\varphi)$ is therefore a \emph{pointed presheaf}---an
object of the slice category $\PSh(C)/\mathbf{y}(c)$, where
$\mathbf{y}$ is the Yoneda embedding.

But a pointed presheaf lives inside a much richer world: the topos
$\Sh(C,J)$ of sheaves on the site.  This topos has a subobject
classifier~$\Omega$, an internal Heyting algebra, exponential objects,
a sheafification functor, geometric morphisms to and from other topoi,
and---crucially---an \emph{internal logic} in which the proposition
$\varphi$ can be interpreted.

None of this structure is visible in the pair $(c,\varphi)$ taken
alone.  Yet all of it is \emph{determined} by the site~$(C,J)$ and
participates in the judgment: the internal logic governs which
propositions are well-formed; the sheaf condition controls how local
data glue; the subobject classifier measures truth values.

\begin{remark}
The passage from $(c,\varphi)$ to a full topos is not a gratuitous
abstraction.  Every topos-theoretic concept we introduce in
Chapter~\ref{ch:stopos} will be exercised in the examples of
Part~\ref{part:examples}.
\end{remark}

\section{The second aspect: $(A,E)$ is a typing judgment}
\label{sec:obs-type}

A carrier type~$A$ together with an evidence term~$E$ satisfying
$\Gamma \vdash E : A$ is a \emph{typing judgment} in the sense of
Martin-L\"of type theory.  It is not merely a pair of a set and an
element.  The derivation witnessing~$\Gamma \vdash E : A$ is a tree
whose internal structure records every logical step: variable lookups,
$\lambda$-abstractions, applications, eliminations.

This derivation tree lives inside a rich proof-theoretic universe:
\begin{itemize}[nosep]
\item Dependent function types $\Pi(x:A).\,B(x)$ and dependent pair
  types $\Sigma(x:A).\,B(x)$.
\item Identity types $\mathrm{Id}_A(a,b)$ with path induction and, in
  the homotopical setting, higher identity types.
\item A universe hierarchy
  $\Type_0 : \Type_1 : \Type_2 : \cdots$.
\item $\beta$-reduction (computation) and $\eta$-expansion
  (extensionality), satisfying subject reduction, confluence, and
  strong normalisation.
\item W-types for general inductive definitions.
\end{itemize}
The pair $(A,E)$, extracted from this universe, discards the context,
the derivation, the reduction theory, and the type formers.  Restoring
them gives the \emph{Proof Structure}~$\PStr$.

\section{The third aspect: $(O,B)$ is a cohomological datum}
\label{sec:obs-cohom}

The obligation cover~$O$ and the obstruction class~$B$ encode the
failure of local sections to glue.  In the language of \v{C}ech
cohomology:
\begin{itemize}[nosep]
\item $O = \{U_i \to X\}$ is an open cover of the site object~$X$.
\item $B \in \check{H}^1(\mathcal{U},\mathcal{F})$ is a first
  cohomology class measuring the obstruction to global gluing.
\end{itemize}
The fundamental theorem of Judgment Geometry states:
\[
  P \models \varphi
  \quad\Longleftrightarrow\quad
  \check{H}^1(\mathcal{U},\mathcal{F}_\varphi) = 0.
\]
But $\check{H}^1$ is only the first layer of a much deeper structure.
The full \v{C}ech complex $C^\bullet(\mathcal{U},\mathcal{F})$ yields
cohomology groups in every degree.  The short exact sequence of sheaves
produces a long exact sequence in cohomology.  The Mayer--Vietoris
sequence relates the cohomology of a union to the cohomology of its
parts.  Spectral sequences connect \v{C}ech cohomology to derived
functor cohomology.  Characteristic classes detect higher-order
twisting.

All of this machinery is implicit in the pair $(O,B)$ but invisible in
the component-wise view.  Making it explicit gives the \emph{Cohomological
Spectrum}~$\CSpec$.

\section{The fourth aspect: $(T,\Pi)$ is a graded derivation}
\label{sec:obs-grade}

The trust grade~$T$ and the provenance trail~$\Pi$ record the quality
of a judgment and the history of how that quality was established.  At
the component level, $T$ ranges over a linearly ordered set
\[
  T_{\mathrm{contra}} \prec T_{\mathrm{unver}} \prec T_{\mathrm{copilot}}
  \prec T_{\mathrm{oracle}} \prec T_{\mathrm{runtime}} \prec
  T_{\mathrm{solver}} \prec T_{\mathrm{proof}},
\]
with operations of conservative join, attenuation, promotion, and
demotion.

This is the shadow of a full algebraic structure.  The grade set~$G$,
equipped with two binary operations $\opG$ (join) and $\otG$ (tensor),
a zero~$\gzero$, and a unit~$\gone$, forms an \emph{idempotent
commutative semiring}---in fact, a \emph{commutative unital quantale}.
The quantale carries:
\begin{itemize}[nosep]
\item A complete lattice structure with $\opG = \bigvee$.
\item A residuation operation $\multimap$ right adjoint to $\otG$,
  making the quantale \emph{closed monoidal}.
\item An evaluation functor from the proof category to the quantale,
  making it a \emph{graded monad}.
\item Stability and convergence conditions governing iterated
  re-evaluation.
\end{itemize}
The pair $(T,\Pi)$ records a single grade and a single trail.
Embedding them in the full quantale-enriched structure gives the
\emph{Graded Evaluation}~$\GEval$.

\section{The thesis: four maximally rich objects}

We have identified four pairs and four ambient structures:
\begin{center}
\renewcommand{\arraystretch}{1.3}
\begin{tabular}{@{}cccc@{}}
\textbf{Pair} & \textbf{Ambient structure} & \textbf{Rich object} &
\textbf{Chapter} \\
\hline
$(c,\varphi)$ & elementary topos with site &
$\STop$ & \ref{ch:stopos} \\
$(A,E)$ & proof-relevant type theory &
$\PStr$ & \ref{ch:pstr} \\
$(O,B)$ & derived \v{C}ech cohomology &
$\CSpec$ & \ref{ch:cspec} \\
$(T,\Pi)$ & quantale-enriched evaluation &
$\GEval$ & \ref{ch:geval}
\end{tabular}
\end{center}

A judgment is a compatible quadruple $(\STop,\PStr,\CSpec,\GEval)$.
The compatibility conditions, formulated as a fibration structure on the
total space of judgments, are the subject of
Chapter~\ref{ch:fibration}.

\section{Why ``rich''}

Each of the four objects is \emph{maximally rich} in a precise sense: it
carries all the internal structure that its ambient mathematical
theory provides.  Concretely:

\begin{enumerate}[label=(\roman*),nosep]
\item \textbf{Internal coherence.}  Each object satisfies the axioms of
  its ambient category (topos axioms, type-theoretic rules, cochain
  complex identities, semiring axioms) without external enforcement.
\item \textbf{Own morphisms.}  Each object lives in a category with a
  well-defined notion of morphism: geometric morphisms for~$\STop$,
  context morphisms for~$\PStr$, chain maps for~$\CSpec$,
  grade-preserving maps for~$\GEval$.
\item \textbf{Own logic.}  Each object carries an internal logic:
  the Mitchell--B\'enabou language of the topos, the type-theoretic
  derivation calculus, the long exact sequence in cohomology, the
  residuated implication $\multimap$ of the quantale.
\end{enumerate}

The component-wise view retains only the \emph{data} of each pair; the
4-tuple retains both the data and the \emph{ambient structure}.  Every
theorem can be stated in the 4-tuple language, and new
theorems---invisible from the flattened component
perspective---become accessible.

\begin{remark}[Enrichment is canonical]
\label{rem:canonical}
Given a set of components $(c,\varphi)$, there is a smallest topos
containing them; given $(A,E)$, a minimal type theory; and so on.
This canonicity is made precise in
Chapter~\ref{ch:equivalence}, where we exhibit an adjoint equivalence
between the structured and flattened views.
\end{remark}

\chapter{The Semantic Topos \texorpdfstring{$\STop$}{S}}
\label{ch:stopos}

The first of our four objects encodes \emph{where} a judgment lives and
\emph{what} it asserts.  At the component level, this information
appears as the pair $(c,\varphi)$---a coordinate and a presheaf
section.  We now develop it into a full elementary topos equipped with
a distinguished presheaf and the complete internal logical apparatus
that the topos provides.

\section{Sites and Grothendieck topologies}
\label{sec:sites}

\begin{definition}[Site]
\label{def:site}
A \emph{site} is a pair $(C,J)$ where $C$ is a small category and
$J$ assigns to each object $U \in \ob(C)$ a collection $J(U)$ of
\emph{covering sieves} on~$U$---families of morphisms with codomain
$U$ closed under precomposition---satisfying:
\begin{enumerate}[label=(\alph*),nosep]
\item \textbf{Maximality.}  The maximal sieve
  $\{f : V \to U \mid V \in \ob(C)\}$ is in $J(U)$.
\item \textbf{Stability.}  If $S \in J(U)$ and $g : V \to U$ is any
  morphism, then the pullback sieve
  $g^*(S) = \{h : W \to V \mid g \circ h \in S\}$ is in $J(V)$.
\item \textbf{Transitivity.}  If $S \in J(U)$ and $R$ is a sieve on
  $U$ such that for every $f : V \to U$ in~$S$ the pullback
  $f^*(R) \in J(V)$, then $R \in J(U)$.
\end{enumerate}
\end{definition}

\begin{example}[Program site]
\label{ex:program-site}
Let $P$ be a program.  The \emph{semantic site} $\mathcal{S}_P$ has:
\begin{itemize}[nosep]
\item Objects: semantic coordinates---modules, classes, functions,
  branches, expressions.
\item Morphisms: scope inclusions $f \hookrightarrow g$ when $f$ is
  defined in the scope of~$g$.
\item Topology: a family $\{U_i \to U\}$ covers~$U$ if the $U_i$
  jointly exhaust every execution path through~$U$.
\end{itemize}
\end{example}

\begin{example}[Discourse site]
\label{ex:discourse-site}
A natural-language discourse~$D$ determines a site whose objects are
discourse segments (sentences, clauses, referring expressions) and
whose coverings reflect anaphoric accessibility: $\{s_i\}$ covers a
discourse referent~$x$ if the $s_i$ jointly determine the reference
of~$x$.
\end{example}

\section{Presheaves and the Yoneda embedding}

\begin{definition}[Presheaf category]
For a small category~$C$, the \emph{presheaf category}
$\PSh(C) = [C^{\op}, \Set]$ has:
\begin{itemize}[nosep]
\item Objects: functors $\mathcal{F} : C^{\op} \to \Set$.
\item Morphisms: natural transformations.
\end{itemize}
\end{definition}

The Yoneda embedding $\mathbf{y} : C \hookrightarrow \PSh(C)$ sends
$U \mapsto \Hom_C(-,U)$ and is fully faithful.  The Yoneda lemma gives
$\Hom_{\PSh(C)}(\mathbf{y}(U), \mathcal{F}) \cong \mathcal{F}(U)$.

\begin{definition}[The datum $\varphi$]
\label{def:phi}
The \emph{judgment presheaf} is a presheaf
$\varphi \in \ob(\PSh(C))$.  For each coordinate
$U \in \ob(C)$, the set $\varphi(U)$ records the local data of
the proposition being judged at stage~$U$.  For each inclusion
$f : V \to U$, the restriction $\varphi(f) : \varphi(U) \to
\varphi(V)$ specialises data from a larger to a smaller scope.
\end{definition}

\section{The sheaf condition and sheafification}
\label{sec:sheafification}

\begin{definition}[Sheaf]
A presheaf $\mathcal{F} \in \PSh(C)$ is a \emph{sheaf} for the
topology~$J$ if, for every covering sieve $S \in J(U)$, the natural
map
\[
  \mathcal{F}(U) \;\longrightarrow\;
  \lim_{(V \to U) \in S}\, \mathcal{F}(V)
\]
is an isomorphism.  Equivalently, every compatible family of local
sections has a unique amalgamation.
\end{definition}

\begin{construction}[Plus construction]
\label{con:plus}
For a presheaf $\mathcal{F}$, the \emph{plus construction} yields
$\mathcal{F}^+$ with
\[
  \mathcal{F}^+(U) \;=\;
  \varinjlim_{S \in J(U)}\;
  \mathrm{Match}(S, \mathcal{F}),
\]
where $\mathrm{Match}(S,\mathcal{F})$ is the set of matching families
for the sieve~$S$.  Iterating twice---$\mathcal{F}^{++}$---yields the
\emph{associated sheaf} (sheafification), and the canonical map
$\mathcal{F} \to \mathcal{F}^{++}$ is universal.
\end{construction}

\begin{definition}[Sheafification adjunction]
\label{def:sheafification}
Let $i : \Sh(C,J) \hookrightarrow \PSh(C)$ be the inclusion.  The
\emph{sheafification functor} $a : \PSh(C) \to \Sh(C,J)$ is the left
adjoint of~$i$:
\[
  a \;\dashv\; i,
  \qquad
  \Hom_{\Sh}(a\mathcal{F},\, \mathcal{G})
  \;\cong\;
  \Hom_{\PSh}(\mathcal{F},\, i\mathcal{G}).
\]
The inclusion~$i$ is fully faithful, so the counit
$a \circ i \xrightarrow{\;\sim\;} \id_{\Sh}$ is an isomorphism.
\end{definition}

\section{The subobject classifier and internal logic}
\label{sec:omega}

\begin{definition}[Subobject classifier]
\label{def:subobj}
A \emph{subobject classifier} in a topos~$\mathcal{E}$ is a monomorphism
$\mathrm{true} : 1 \to \Omega$ such that for every monomorphism
$m : S \hookrightarrow X$ there is a unique \emph{characteristic map}
$\chi_m : X \to \Omega$ making the square
\[
\begin{tikzcd}
  S \arrow[r] \arrow[d, hook, "m"'] & 1 \arrow[d, "\mathrm{true}"] \\
  X \arrow[r, "\chi_m"'] & \Omega
\end{tikzcd}
\]
a pullback.
\end{definition}

\begin{theorem}
\label{thm:omega-exists}
In the sheaf topos $\Sh(C,J)$, the subobject classifier exists and is
given by
\[
  \Omega(U) \;=\; \{\text{closed sieves on } U\},
\]
where a sieve~$S$ on~$U$ is \emph{$J$-closed} if, whenever
$f^*(S) \in J(V)$, then $f \in S$.  The map
$\mathrm{true}_U : \{*\} \to \Omega(U)$ picks the maximal sieve.
\end{theorem}

\begin{proposition}[Heyting algebra structure]
\label{prop:heyting}
For each $U \in \ob(C)$, the set $\Omega(U)$ carries a Heyting algebra
structure:
\begin{align*}
  (S_1 \wedge S_2)(U) &= S_1 \cap S_2, \\
  (S_1 \vee S_2)(U) &= j(S_1 \cup S_2), \\
  (S_1 \Rightarrow S_2)(U) &= \{f : V \to U \mid
    \text{for all } g : W \to V,\;
    f \circ g \in S_1 \;\Longrightarrow\;
    f \circ g \in S_2 \},
\end{align*}
where $j$ denotes the closure operation associated to~$J$.  In
particular, $\Omega$ is in general \emph{not} Boolean:
$S \vee \neg S$ need not equal the maximal sieve.
\end{proposition}

\begin{definition}[Internal quantifiers]
\label{def:quantifiers}
For each morphism $f : Y \to X$ in the topos, the pullback functor
$f^* : \mathrm{Sub}(X) \to \mathrm{Sub}(Y)$ on subobject lattices has:
\begin{itemize}[nosep]
\item a left adjoint $\exists_f \dashv f^*$ ---
  the \emph{internal existential} along~$f$;
\item a right adjoint $f^* \dashv \forall_f$ ---
  the \emph{internal universal} along~$f$.
\end{itemize}
These satisfy the Beck--Chevalley condition: for every pullback square,
$\exists$ and $\forall$ commute with pullback.
\end{definition}

\begin{remark}[The Mitchell--B\'enabou language]
\label{rem:mitchell-benabou}
The internal logic of a topos~$\mathcal{E}$ is formalised by the
\emph{Mitchell--B\'enabou language}: terms are generalised elements
(morphisms into objects of~$\mathcal{E}$), formulas are morphisms into
$\Omega$, and the logical connectives $\wedge, \vee, \Rightarrow,
\neg, \forall, \exists$ are the natural transformations described
above.  The semantics is Kripke--Joyal: a formula
$\varphi$ \emph{holds at stage~$U$} if the corresponding sieve is
$J$-dense at~$U$.  This means the topos provides a
\emph{constructive higher-order logic} in which the judgment
presheaf~$\varphi$ can be interpreted without any external logical
framework.
\end{remark}

\section{Internal equality and Lawvere--Tierney topologies}

\begin{definition}[Internal equality]
For each object $X$ in the topos, the \emph{equality predicate}
$=_X : X \times X \to \Omega$ is the characteristic map of the diagonal
$\Delta_X : X \hookrightarrow X \times X$.
\end{definition}

\begin{definition}[Lawvere--Tierney topology]
\label{def:lt-topology}
A \emph{Lawvere--Tierney topology} (or \emph{local operator}) on a
topos~$\mathcal{E}$ is an endomorphism $j : \Omega \to \Omega$
satisfying:
\begin{enumerate}[label=(\alph*),nosep]
\item $j \circ \mathrm{true} = \mathrm{true}$,
\item $j \circ j = j$,
\item $j \circ \wedge = \wedge \circ (j \times j)$.
\end{enumerate}
\end{definition}

\begin{proposition}
Every Lawvere--Tierney topology~$j$ determines a sub-topos
$\Sh_j(\mathcal{E}) \hookrightarrow \mathcal{E}$ whose objects are the
$j$-sheaves, and conversely every sub-topos arises from a unique~$j$.
The operator~$j$ provides a \emph{modal operator} on truth values:
$j(\varphi)$ is the $j$-closure of~$\varphi$, interpretable as
``$\varphi$ holds $j$-locally.''
\end{proposition}

In Judgment Geometry, different Lawvere--Tierney topologies correspond to
different notions of \emph{local truth}---e.g., ``true on all
branches,'' ``true up to syntactic $\alpha$-equivalence,'' or ``true
modulo discourse context.''

\section{Points and stalks}

\begin{definition}[Point of a topos]
\label{def:point}
A \emph{point} of $\Sh(C,J)$ is a geometric morphism
$p : \Set \to \Sh(C,J)$, i.e., an adjoint pair
$p^* \dashv p_*$ with $p^*$ left exact (preserving finite limits).
The \emph{stalk} of a sheaf~$\mathcal{F}$ at~$p$ is the set
$\mathcal{F}_p = p^*(\mathcal{F})$.
\end{definition}

\begin{theorem}[Deligne]
\label{thm:deligne}
If the site $(C,J)$ is \emph{coherent}---every covering sieve
contains a finite sub-cover---then the sheaf topos $\Sh(C,J)$ has
\emph{enough points}: a morphism of sheaves is an isomorphism if and
only if it induces isomorphisms on all stalks.
\end{theorem}

\begin{remark}
Program sites are typically coherent (functions have finitely many
branches), so Deligne's theorem applies: local-to-global arguments via
stalks are valid.  Discourse sites require more care; infinite
anaphoric chains may violate coherence.
\end{remark}

\section{Slice topoi and dependent types}

\begin{theorem}
\label{thm:slice-topos}
For any object $X$ of a topos~$\mathcal{E}$, the slice category
$\mathcal{E}/X$ is again a topos.  Its subobject classifier is
$\Omega_X = \Omega \times X \to X$ (the projection), and the
forgetful functor $\Sigma_X : \mathcal{E}/X \to \mathcal{E}$ has
both a left adjoint $X^* = (-) \times X$ and a right adjoint
$\Pi_X = (-)^X$ (internal hom over~$X$).
\end{theorem}

In Judgment Geometry, the slice topos $\STop/\varphi$ encodes
\emph{dependent judgments}: judgments that depend on the data
of~$\varphi$.  An object of $\STop/\varphi$ is a sheaf $\mathcal{G}$
together with a map $\mathcal{G} \to \varphi$, representing a family of
types indexed by the local sections of~$\varphi$.

\section{Geometric morphisms}
\label{sec:geom-morph}

\begin{definition}[Geometric morphism]
\label{def:geom-morph}
A \emph{geometric morphism} $f : \mathcal{E} \to \mathcal{F}$ between
topoi is an adjoint pair $f^* \dashv f_*$ where the \emph{inverse
image} $f^*$ preserves finite limits.  A geometric morphism is:
\begin{itemize}[nosep]
\item \emph{essential} if $f^*$ has a further left adjoint $f_!$,
  giving $f_! \dashv f^* \dashv f_*$;
\item \emph{surjective} (or a \emph{surjection}) if $f^*$ is faithful;
\item \emph{an inclusion} (or \emph{embedding}) if $f_*$ is fully
  faithful.
\end{itemize}
\end{definition}

\begin{theorem}[Hyperconnected--localic factorisation]
\label{thm:factorisation}
Every geometric morphism $f : \mathcal{E} \to \mathcal{F}$ factors
uniquely as
\[
  \mathcal{E} \xrightarrow{\;f_h\;} \mathcal{G}
  \xrightarrow{\;f_l\;} \mathcal{F},
\]
where $f_h$ is hyperconnected (connected and locally connected) and
$f_l$ is localic (the inverse image is a Heyting algebra homomorphism
on subobject lattices).
\end{theorem}

In practice, geometric morphisms encode \emph{change of site}: if the
program is refactored (renaming modules, splitting functions), the
old and new sites are related by a geometric morphism, and the
factorisation theorem decomposes the refactoring into a ``topological''
part (the localic factor, changing the underlying space) and a
``logical'' part (the hyperconnected factor, changing the fibres).

\section{The complete definition}

We now assemble the full definition of the Semantic Topos.

\begin{definition}[Semantic Topos $\STop$]
\label{def:stop}
A \emph{Semantic Topos} is a tuple
\[
  \STop \;=\;
  \bigl(C,\; J,\; \varphi,\; \Omega,\; \mathrm{Sub},\;
  \wedge,\, \vee,\, \Rightarrow,\,
  \forall,\, \exists,\; {=},\;
  a,\; \mathrm{Pts},\; \tau \bigr)
\]
where:
\begin{enumerate}[label=(\roman*),nosep]
\item $(C,J)$ is a site (Definition~\ref{def:site});
\item $\varphi \in \ob(\PSh(C))$ is the judgment presheaf
  (Definition~\ref{def:phi});
\item $\Omega$ is the subobject classifier of $\Sh(C,J)$
  (Theorem~\ref{thm:omega-exists});
\item $\mathrm{Sub}$ is the subobject functor,
  $\wedge, \vee, \Rightarrow$ are the Heyting operations on~$\Omega$
  (Proposition~\ref{prop:heyting});
\item $\forall, \exists$ are the internal quantifiers
  (Definition~\ref{def:quantifiers});
\item ${=}$ is internal equality via the diagonal;
\item $a : \PSh(C) \to \Sh(C,J)$ is sheafification
  (Definition~\ref{def:sheafification});
\item $\mathrm{Pts}$ is the class of points of $\Sh(C,J)$
  (Definition~\ref{def:point});
\item $\tau$ is a choice of Lawvere--Tierney topology
  $j : \Omega \to \Omega$ (Definition~\ref{def:lt-topology}),
  determining the operative notion of local truth.
\end{enumerate}
\end{definition}

\begin{remark}
The data (iii)--(viii) are \emph{determined} by the site $(C,J)$; they
are not additional choices.  We list them to emphasise that $\STop$
\emph{carries} all this structure, making it available for the internal
logic of judgments.  Only~$\varphi$ (the datum being judged) and
$\tau$ (the chosen notion of local truth) involve genuine choices.
\end{remark}

\begin{definition}[Morphism of Semantic Topoi]
\label{def:stop-morph}
A \emph{morphism} $f : \STop_1 \to \STop_2$ is a geometric morphism
$f = (f^*,f_*) : \Sh(C_1,J_1) \to \Sh(C_2,J_2)$
(Definition~\ref{def:geom-morph}) together with a natural isomorphism
$\alpha : f^*(\varphi_2) \xrightarrow{\;\sim\;} \varphi_1$ that is
compatible with the Lawvere--Tierney topologies:
$f^* \circ j_2 = j_1 \circ f^*$ on~$\Omega_1$.
\end{definition}

We write $\mathbf{STop}$ for the category of Semantic Topoi and their
morphisms.

\section{Examples}
\label{sec:stop-examples}

\begin{example}[Program AST topos]
\label{ex:ast-topos}
For a typed functional program, the site objects are AST nodes
(module, function, let-binding, pattern-match branch, expression),
morphisms are containment, and covering families are exhaustive case
splits.  The presheaf $\varphi$ assigns to each node the type
specification it must satisfy.  The subobject classifier $\Omega$
assigns to each node the lattice of ``locally valid partial
specifications.''  Sheafification forces local specifications to be
globally consistent.
\end{example}

\begin{example}[Natural-language discourse topos]
\label{ex:discourse-topos}
For a discourse, the site objects are Discourse Representation
Structures (DRSs) in the sense of Kamp--Reyle, morphisms are
accessibility relations, and covers are exhaustive sets of accessible
antecedents for a pronoun.  The presheaf $\varphi$ assigns to each DRS
the set of possible referents.  The internal logic is intuitionistic:
$\varphi \vee \neg\varphi$ fails when a pronoun is genuinely ambiguous.
\end{example}

\begin{example}[Poetic form topos]
\label{ex:poetry-topos}
For a poem, site objects are formal positions (line, foot, syllable);
covers require that metrical constraints be checked on all feet within
a line.  The presheaf $\varphi$ assigns to each position the set of
admissible stress patterns.  Sheafification forces local metrical
constraints to be globally consistent, detecting enjambment and
intentional metrical violations.
\end{example}

\begin{example}[Mathematical theory topos]
For a mathematical theory~$T$, the site is the syntactic category
$\mathcal{C}_T$ whose objects are formulas-in-context and whose
morphisms are provable entailments.  The topology is generated by
finite disjunctions.  The presheaf $\varphi$ is the definable set
corresponding to the formula under scrutiny.  Sheafification yields the
classifying topos $\Set[T]$, and geometric morphisms
$\Set \to \Set[T]$ are precisely the models of~$T$.
\end{example}

\chapter{The Proof Structure \texorpdfstring{$\PStr$}{P}}
\label{ch:pstr}

The second object encodes the \emph{evidence} supporting a judgment.
At the component level, the pair $(A,E)$---a type and a term inhabiting
it---records the bare fact that evidence exists.  We now develop the
full proof-relevant type-theoretic universe in which this pair lives:
dependent types, derivation trees, reduction theory,
identity types, and inductive definitions.

\section{Contexts and well-formedness}
\label{sec:contexts}

\begin{definition}[Context]
\label{def:context}
A \emph{context} is a telescope of typed variable declarations
\[
  \Gamma \;=\; (x_1 : A_1,\; x_2 : A_2(x_1),\;
  \ldots,\; x_n : A_n(x_1,\ldots,x_{n-1})),
\]
where each $A_k$ may depend on the preceding variables.  The empty
context is written $\cdot$.
\end{definition}

\begin{definition}[Context well-formedness]
Well-formedness $\vdash \Gamma \;\mathsf{ctx}$ is defined inductively:
\begin{enumerate}[label=(\alph*),nosep]
\item $\vdash \cdot \;\mathsf{ctx}$.
\item If $\vdash \Gamma \;\mathsf{ctx}$ and
  $\Gamma \vdash A : \Type_i$ for some~$i$, then
  $\vdash (\Gamma, x : A) \;\mathsf{ctx}$, provided
  $x \notin \mathrm{dom}(\Gamma)$.
\end{enumerate}
\end{definition}

Standard structural rules govern the manipulation of contexts:

\begin{definition}[Structural rules]
\label{def:structural}
Let $\Gamma$ be a well-formed context.
\begin{itemize}[nosep]
\item \textbf{Weakening.} If $\Gamma \vdash \mathcal{J}$ and
  $\vdash (\Gamma, x:A) \;\mathsf{ctx}$, then
  $\Gamma, x:A \vdash \mathcal{J}$.
\item \textbf{Contraction.} If $\Gamma, x:A, y:A \vdash \mathcal{J}$,
  then $\Gamma, x:A \vdash \mathcal{J}[y := x]$.
\item \textbf{Exchange.} If $x$ does not occur in~$B$ and $y$ does not
  occur in~$A$, then $\Gamma, x:A, y:B \vdash \mathcal{J}$ implies
  $\Gamma, y:B, x:A \vdash \mathcal{J}$.
\end{itemize}
\end{definition}

\section{Types and universe hierarchy}
\label{sec:universes}

\begin{definition}[Universe hierarchy]
\label{def:universe}
The \emph{universe hierarchy} is a countable tower
\[
  \Type_0 \;:\; \Type_1 \;:\; \Type_2 \;:\; \cdots
\]
with a cumulativity rule: if $\Gamma \vdash A : \Type_i$ then
$\Gamma \vdash A : \Type_{i+1}$.  We adopt the Russell convention:
a type $A : \Type_i$ is simultaneously a term of $\Type_i$ and a
type whose terms $a : A$ can be formed.
\end{definition}

\begin{remark}[Universe polymorphism]
In practice, we work \emph{universe-polymorphically}: definitions and
theorems are quantified over universe levels, avoiding the need to fix
a specific $\Type_i$.  This is essential when the type~$A$ in a
judgment may itself be a universe (e.g., $A = \Type_0$).
\end{remark}

\section{Evidence terms and the $\lambda$-calculus}

\begin{definition}[Evidence terms]
\label{def:terms}
The \emph{evidence terms} are generated by the grammar
\[
  E \;::=\; x \;\mid\; \lambda x.\, E \;\mid\; E_1\; E_2
  \;\mid\; (E_1, E_2) \;\mid\; \pi_1\, E \;\mid\; \pi_2\, E
  \;\mid\; \mathrm{let}\; x = E_1 \;\mathrm{in}\; E_2
  \;\mid\; \cdots
\]
where $x$ ranges over variables and ``$\cdots$'' stands for the
constructors and eliminators of the specific type formers (identity
types, W-types, etc.\ introduced below).  Each term is assigned a type
by a derivation of the form $\Gamma \vdash E : A$.
\end{definition}

\begin{definition}[Explicit substitution]
The \emph{substitution} $E[x := N]$ replaces every free occurrence
of~$x$ in~$E$ by~$N$, with the standard capture-avoiding convention.
Substitution satisfies:
\begin{enumerate}[label=(\alph*),nosep]
\item $x[x := N] = N$.
\item $y[x := N] = y$ if $y \neq x$.
\item $(\lambda y.\, E)[x := N] = \lambda y.\, E[x := N]$ if
  $y \neq x$ and $y \notin \mathrm{FV}(N)$ (rename otherwise).
\item $(E_1 \; E_2)[x := N] = E_1[x := N] \; E_2[x := N]$.
\end{enumerate}
\end{definition}

\section{Typing derivations}
\label{sec:derivations}

\begin{definition}[Typing derivation]
\label{def:derivation}
A \emph{typing derivation} $\mathcal{D}$ is a finite tree whose nodes
are typing judgments and whose edges connect premises to conclusions
according to the inference rules of the type theory.  Each node is
labelled by the name of the rule applied.

Formally, $\mathcal{D}$ is an element of the inductively defined set
of derivation trees:
\[
  \mathcal{D} \;=\;
  \frac{\mathcal{D}_1 \quad \cdots \quad \mathcal{D}_k}
  {\Gamma \vdash E : A}\;\; (\text{rule-name}),
\]
where $\mathcal{D}_1, \ldots, \mathcal{D}_k$ are the sub-derivations
for the premises.
\end{definition}

\begin{remark}[Proof relevance]
\label{rem:proof-relevance}
Two terms $E$ and $E'$ may both inhabit type~$A$ in context~$\Gamma$
but have \emph{different} derivation trees.  The derivation is
part of the judgment's evidence: it records not just \emph{that} $A$ is
inhabited, but \emph{how}.  This is the essence of
\emph{proof relevance}: the proof matters, not just its existence.
In Judgment Geometry, the derivation tree~$\mathcal{D}$ is the
structural skeleton of the evidence.
\end{remark}

\section{Reduction theory}
\label{sec:reduction}

\begin{definition}[$\beta$-reduction]
\label{def:beta}
The \emph{$\beta$-reduction} relation $\to_\beta$ is generated by:
\[
  (\lambda x.\, M)\; N \;\to_\beta\; M[x := N].
\]
This is the fundamental \emph{computation} rule: applying a function
to an argument computes by substitution.
\end{definition}

\begin{definition}[$\eta$-expansion]
\label{def:eta}
For function types, the \emph{$\eta$-expansion} rule is:
\[
  f \;\to_\eta\; \lambda x.\, f\; x
  \qquad (x \notin \mathrm{FV}(f)).
\]
This enforces \emph{extensionality}: two functions that agree on all
inputs are equal.  Analogous $\eta$-rules exist for pair types
($p \to_\eta (\pi_1\, p,\, \pi_2\, p)$) and other type formers.
\end{definition}

\begin{theorem}[Subject reduction]
\label{thm:subject-reduction}
If $\Gamma \vdash E : A$ and $E \to_\beta E'$, then
$\Gamma \vdash E' : A$.  That is, $\beta$-reduction preserves typing.
\end{theorem}

\begin{theorem}[Confluence]
\label{thm:confluence}
The $\beta$-reduction relation is confluent (Church--Rosser): if
$E \to_\beta^* E_1$ and $E \to_\beta^* E_2$, then there exists~$E_3$
with $E_1 \to_\beta^* E_3$ and $E_2 \to_\beta^* E_3$.
\end{theorem}

\begin{theorem}[Strong normalisation]
\label{thm:sn}
Every well-typed term in the type-theoretic fragment (without general
recursion) is strongly normalising: every reduction sequence from~$E$
terminates.
\end{theorem}

\section{Dependent type formers}
\label{sec:type-formers}

\begin{definition}[Dependent product]
\label{def:pi}
Given $\Gamma \vdash A : \Type_i$ and
$\Gamma, x : A \vdash B(x) : \Type_j$, the \emph{dependent product}
$\Pi(x : A).\, B(x) : \Type_{\max(i,j)}$ has:
\begin{itemize}[nosep]
\item \textbf{Introduction:}
  $\Gamma \vdash \lambda x.\, e : \Pi(x:A).\, B(x)$
  when $\Gamma, x : A \vdash e : B(x)$.
\item \textbf{Elimination:}
  $\Gamma \vdash f\; a : B(a)$
  when $\Gamma \vdash f : \Pi(x:A).\, B(x)$ and $\Gamma \vdash a : A$.
\item \textbf{Computation:}
  $(\lambda x.\, e)\; a \;=_\beta\; e[x := a]$.
\item \textbf{Uniqueness:}
  $f \;=_\eta\; \lambda x.\, f\; x$.
\end{itemize}
When $B$ does not depend on~$x$, we recover the simple function type
$A \to B$.
\end{definition}

\begin{definition}[Dependent sum]
\label{def:sigma}
Given $\Gamma \vdash A : \Type_i$ and
$\Gamma, x : A \vdash B(x) : \Type_j$, the \emph{dependent sum}
$\Sigma(x : A).\, B(x) : \Type_{\max(i,j)}$ has:
\begin{itemize}[nosep]
\item \textbf{Introduction:}
  $\Gamma \vdash (a, b) : \Sigma(x:A).\, B(x)$
  when $\Gamma \vdash a : A$ and $\Gamma \vdash b : B(a)$.
\item \textbf{Elimination:}
  $\pi_1 : \Sigma(x:A).\, B(x) \to A$ and
  $\pi_2 : (p : \Sigma(x:A).\, B(x)) \to B(\pi_1\, p)$.
\item \textbf{Computation:}
  $\pi_1(a,b) =_\beta a$ and $\pi_2(a,b) =_\beta b$.
\item \textbf{Uniqueness:}
  $p =_\eta (\pi_1\, p,\, \pi_2\, p)$.
\end{itemize}
When $B$ does not depend on~$x$, we recover the product type
$A \times B$.
\end{definition}

\section{Identity types}
\label{sec:identity}

\begin{definition}[Identity type]
\label{def:identity}
For $\Gamma \vdash a, b : A$, the \emph{identity type}
$\mathrm{Id}_A(a,b) : \Type_i$ (where $A : \Type_i$) has:
\begin{itemize}[nosep]
\item \textbf{Introduction:}
  $\mathrm{refl}_a : \mathrm{Id}_A(a,a)$ --- the reflexivity witness.
\item \textbf{Elimination (J-rule):}  Given a type family
  $C : \Pi(x,y:A).\, \mathrm{Id}_A(x,y) \to \Type_j$ and a term
  $d : \Pi(x:A).\, C(x,x,\mathrm{refl}_x)$, there is a term
  \[
    \mathrm{J}(C,d,a,b,p) : C(a,b,p)
  \]
  for every $p : \mathrm{Id}_A(a,b)$, with computation rule
  $\mathrm{J}(C,d,a,a,\mathrm{refl}_a) =_\beta d(a)$.
\end{itemize}
\end{definition}

\begin{remark}[Higher identity types]
The identity type can be iterated: given $p, q : \mathrm{Id}_A(a,b)$,
we can form $\mathrm{Id}_{\mathrm{Id}_A(a,b)}(p,q)$---a type whose
inhabitants are \emph{paths between paths}.  This gives each type~$A$
the structure of an $\infty$-groupoid, connecting type theory to
homotopy theory.
\end{remark}

\begin{remark}[Univalence]
The \emph{univalence axiom} states that the canonical map
$(A =_{\Type} B) \to (A \simeq B)$, sending a path to the induced
equivalence, is itself an equivalence.  Under univalence, equivalent
types are identical as elements of the universe.  We do not require
univalence in the base theory but note that it is compatible with all
our constructions.
\end{remark}

\section{W-types and general induction}
\label{sec:wtypes}

\begin{definition}[W-type]
\label{def:wtype}
Given $\Gamma \vdash A : \Type_i$ and
$\Gamma, x : A \vdash B(x) : \Type_i$, the \emph{W-type}
$W(x : A).\, B(x) : \Type_i$ is the type of well-founded trees with:
\begin{itemize}[nosep]
\item \textbf{Constructor:}
  $\mathrm{sup} : \Pi(a : A).\, (B(a) \to W(x:A).\, B(x))
  \to W(x:A).\, B(x)$.
  Intuitively: a tree consists of a label $a : A$ at the root and a
  $B(a)$-indexed family of subtrees.
\item \textbf{Eliminator:}
  $\mathrm{wrec}$ --- the dependent recursor, which maps out of the
  W-type by providing, for each $a : A$ and each family of subtrees
  $f : B(a) \to W$, a value in the target type, given recursive access
  to the values on subtrees.
\item \textbf{Computation:}
  $\mathrm{wrec}(h, \mathrm{sup}(a,f))
  =_\beta h(a, f, \lambda b.\, \mathrm{wrec}(h, f(b)))$.
\end{itemize}
\end{definition}

\begin{proposition}
W-types suffice to encode all finitary inductive types: natural
numbers ($A = \mathbf{2}$, $B(\mathtt{zero}) = \mathbf{0}$,
$B(\mathtt{succ}) = \mathbf{1}$), lists, ordinal notations, and the
syntax trees of typed $\lambda$-calculi.
\end{proposition}

\section{The proof category}
\label{sec:proof-category}

\begin{construction}[Proof category]
\label{con:proof-category}
Given a context~$\Gamma$, we construct the \emph{proof category}
$\mathbf{Proof}_\Gamma$:
\begin{itemize}[nosep]
\item \textbf{Objects:} types $A$ with $\Gamma \vdash A : \Type_i$
  for some~$i$.
\item \textbf{Morphisms:} $\Hom(A,B) = \{E \mid
  \Gamma \vdash E : A \to B\} / {=_{\beta\eta}}$---terms of function
  type modulo $\beta\eta$-equivalence.
\item \textbf{Composition:} $g \circ f = \lambda x.\, g\;(f\;x)$.
\item \textbf{Identity:} $\id_A = \lambda x.\, x$.
\end{itemize}
\end{construction}

\begin{theorem}
\label{thm:lccc}
The proof category $\mathbf{Proof}_\Gamma$ is a locally cartesian
closed category (LCCC):
\begin{enumerate}[label=(\alph*),nosep]
\item It has a terminal object~$\mathbf{1}$ (the unit type).
\item It has pullbacks (constructed via dependent sums and identity
  types).
\item For every morphism $f : A \to B$, the pullback functor
  $f^* : \mathbf{Proof}_\Gamma / B \to
  \mathbf{Proof}_\Gamma / A$ has a right adjoint $\Pi_f$
  (dependent product along~$f$) and a left adjoint $\Sigma_f$
  (dependent sum along~$f$).
\end{enumerate}
Moreover, there is an equivalence of categories between LCCCs with
additional structure (universes, W-types) and dependent type theories
with the corresponding type formers.
\end{theorem}

\section{The complete definition}

\begin{definition}[Proof Structure $\PStr$]
\label{def:pstr}
A \emph{Proof Structure} is a tuple
\[
  \PStr \;=\;
  (\Gamma,\; A,\; E,\; \mathcal{D},\; {\to_\beta},\;
  {\to_\eta},\; \mathcal{U},\;
  \mathrm{Id},\; \Sigma,\; \Pi,\; W,\; \mathrm{Elim})
\]
where:
\begin{enumerate}[label=(\roman*),nosep]
\item $\Gamma$ is a well-formed context (Definition~\ref{def:context});
\item $A$ is a type with $\Gamma \vdash A : \Type_i$ for some~$i$;
\item $E$ is an evidence term with $\Gamma \vdash E : A$
  (Definition~\ref{def:terms});
\item $\mathcal{D}$ is a derivation tree witnessing
  $\Gamma \vdash E : A$ (Definition~\ref{def:derivation});
\item $\to_\beta$ is the $\beta$-reduction relation
  (Definition~\ref{def:beta});
\item $\to_\eta$ is the $\eta$-expansion relation
  (Definition~\ref{def:eta});
\item $\mathcal{U} = (\Type_i)_{i \in \mathbb{N}}$ is the universe
  hierarchy (Definition~\ref{def:universe});
\item $\mathrm{Id}$ denotes the identity type former
  (Definition~\ref{def:identity});
\item $\Sigma$ denotes the dependent sum type former
  (Definition~\ref{def:sigma});
\item $\Pi$ denotes the dependent product type former
  (Definition~\ref{def:pi});
\item $W$ denotes the W-type former
  (Definition~\ref{def:wtype});
\item $\mathrm{Elim}$ is the collection of eliminators (induction
  principles) for all type formers.
\end{enumerate}
\end{definition}

\begin{definition}[Morphism of Proof Structures]
\label{def:pstr-morph}
A \emph{morphism} $\sigma : \PStr_1 \to \PStr_2$ is a
\emph{context morphism} (substitution)
$\sigma : \Gamma_2 \to \Gamma_1$---a family of terms
$(\sigma(x_k))_{x_k \in \mathrm{dom}(\Gamma_2)}$ with
$\Gamma_1 \vdash \sigma(x_k) : A_k[\sigma]$---such that
$A_2[\sigma] = A_1$ and $E_2[\sigma] =_{\beta\eta} E_1$, and the
derivation~$\mathcal{D}_2$ transports along~$\sigma$ to a derivation
of $\Gamma_1 \vdash E_1 : A_1$.
\end{definition}

We write $\mathbf{PStr}$ for the category of Proof Structures and their
morphisms.

\section{Examples}

\begin{example}[Type-checking derivation]
\label{ex:type-checking}
A type-checker for a dependently typed language, given a term~$E$ and a
type~$A$ in context~$\Gamma$, constructs a Proof Structure: the
derivation tree~$\mathcal{D}$ is the trace of the type-checking
algorithm.  The $\beta$-reduction relation governs normalisation during
type-checking (e.g., reducing $(\lambda x.\,x)\;3$ to~$3$ to
check that $3 : \mathbb{N}$).
\end{example}

\begin{example}[Semantic parse]
\label{ex:semantic-parse}
A Montague-style semantic parse of a sentence produces a $\PStr$:
$\Gamma$ contains typed discourse referents;
$A = \mathrm{Prop}$ (the type of propositions);
$E$ is the $\lambda$-term
representing the sentence's meaning (e.g., $\lambda x.\,
\mathsf{dog}(x) \wedge \mathsf{runs}(x)$);
$\mathcal{D}$ is the derivation tree of the combinatory categorial
grammar.
\end{example}

\begin{example}[Proof of a theorem]
\label{ex:formal-proof}
A formal proof in a proof assistant (Coq, Lean, Agda) is a Proof
Structure.  The context~$\Gamma$ contains the hypotheses; $A$ is the
theorem statement; $E$ is the proof term; $\mathcal{D}$ is the
elaborated derivation; $\to_\beta$ governs reduction during proof
checking.
\end{example}

\chapter{The Cohomological Spectrum \texorpdfstring{$\CSpec$}{H}}
\label{ch:cspec}

The third object encodes the \emph{obstructions} that prevent local
judgments from assembling into a global one.  At the component level,
the pair $(O,B)$ records a cover and a first cohomology class.  We now
develop the full \v{C}ech cohomological apparatus: the cochain
complex in every degree, the long exact sequence, spectral sequences,
and the derived category.

\section{Covers and their combinatorics}
\label{sec:covers}

\begin{definition}[Cover]
\label{def:cover}
Let $X \in \ob(C)$ be an object of the site.  A \emph{cover} of~$X$
is a family $\mathcal{U} = \{U_i \to X\}_{i \in I}$ that generates a
covering sieve in the topology~$J$.  A \emph{refinement}
$\mathcal{V} \to \mathcal{U}$ is a cover $\mathcal{V} = \{V_j \to X\}$
together with a function $\lambda : J' \to I$ and morphisms
$V_j \to U_{\lambda(j)}$ over~$X$.
\end{definition}

\begin{definition}[Nerve]
\label{def:nerve}
The \emph{nerve} $N(\mathcal{U})$ of a cover $\mathcal{U}$ is the
simplicial set with
\[
  N(\mathcal{U})_n \;=\;
  \coprod_{i_0, \ldots, i_n \in I}
  (U_{i_0} \times_X U_{i_1} \times_X \cdots \times_X U_{i_n}),
\]
with face maps given by omitting a factor and degeneracy maps given by
repeating a factor.  When the $U_i$ are open subsets of a space, the
fibre products are intersections:
$U_{i_0} \times_X \cdots \times_X U_{i_n} = U_{i_0} \cap \cdots \cap U_{i_n}$.
\end{definition}

\section{The \v{C}ech complex}
\label{sec:cech-complex}

\begin{definition}[\v{C}ech complex]
\label{def:cech}
Let $\mathcal{F}$ be a presheaf on~$C$ and $\mathcal{U} = \{U_i \to X\}$
a cover.  The \emph{\v{C}ech complex}
$\Cech^\bullet(\mathcal{U}, \mathcal{F})$ is the cochain complex:
\begin{align*}
  \Cech^0(\mathcal{U}, \mathcal{F})
  &= \prod_{i \in I} \mathcal{F}(U_i),
  \\[3pt]
  \Cech^1(\mathcal{U}, \mathcal{F})
  &= \prod_{i < j} \mathcal{F}(U_i \cap U_j),
  \\[3pt]
  \Cech^2(\mathcal{U}, \mathcal{F})
  &= \prod_{i < j < k} \mathcal{F}(U_i \cap U_j \cap U_k),
  \\[3pt]
  \Cech^n(\mathcal{U}, \mathcal{F})
  &= \prod_{i_0 < \cdots < i_n}
  \mathcal{F}(U_{i_0} \cap \cdots \cap U_{i_n}).
\end{align*}
\end{definition}

\begin{definition}[Coboundary maps]
\label{def:coboundary-p2}
The \emph{coboundary} $\delta^n : \Cech^n \to \Cech^{n+1}$ is defined
by alternating restrictions:
\[
  (\delta^n \alpha)_{i_0, \ldots, i_{n+1}}
  \;=\;
  \sum_{k=0}^{n+1} (-1)^k \;
  \alpha_{i_0, \ldots, \widehat{i_k}, \ldots, i_{n+1}}
  \Big|_{U_{i_0} \cap \cdots \cap U_{i_{n+1}}},
\]
where $\widehat{i_k}$ indicates omission. In additive notation, when
$\mathcal{F}$ takes values in abelian groups, the coboundary for $n=0$
is the classical formula
\[
  (\delta^0 s)_{ij} \;=\;
  s_j\big|_{U_i \cap U_j} - s_i\big|_{U_i \cap U_j}.
\]
\end{definition}

\begin{proposition}
\label{prop:d-squared}
$\delta^{n+1} \circ \delta^n = 0$ for all~$n \geq 0$.
\end{proposition}

\begin{proof}
Each term in the double sum
$(\delta^{n+1} \circ \delta^n)(\alpha)$ appears exactly twice with
opposite signs, via the identity
\[
  \sum_{k < l} (-1)^{k+l} + \sum_{l \leq k} (-1)^{k+l+1} \;=\; 0.
  \qedhere
\]
\end{proof}

\section{\v{C}ech cohomology groups}
\label{sec:cech-cohomology}

\begin{definition}[\v{C}ech cohomology]
\label{def:cech-cohom}
The \emph{$n$-th \v{C}ech cohomology group} of $\mathcal{F}$ with
respect to the cover $\mathcal{U}$ is
\[
  \check{H}^n(\mathcal{U}, \mathcal{F})
  \;=\;
  \frac{\ker(\delta^n : \Cech^n \to \Cech^{n+1})}
       {\mathrm{im}(\delta^{n-1} : \Cech^{n-1} \to \Cech^n)},
\]
with the convention $\Cech^{-1} = 0$.
\end{definition}

We now interpret each degree in the context of Judgment Geometry.

\begin{proposition}[Degree-by-degree interpretation]
\label{prop:degree-interp}
\hfill
\begin{enumerate}[label=(\alph*),nosep]
\item $\check{H}^0(\mathcal{U}, \mathcal{F}) = \ker\, \delta^0
  = \Gamma(X, \mathcal{F})$: the space of \emph{global sections},
  i.e., locally defined data that are already globally consistent.
\item $\check{H}^1(\mathcal{U}, \mathcal{F})
  = \ker\, \delta^1 / \mathrm{im}\, \delta^0$: classes of
  \emph{$1$-cocycles modulo coboundaries}.  A non-trivial class is an
  obstruction to gluing---exactly the JuGeo obstruction~$B$.
\item $\check{H}^2(\mathcal{U}, \mathcal{F})$: obstructions to
  \emph{modifying} gluings.  When $H^1 \neq 0$, one may ask whether the
  cocycle can be trivialized by changing the cover; $H^2$ measures
  the secondary obstruction.  This degree illustrates the advantage
  of the full spectrum over isolated components.
\item Higher $\check{H}^n$: higher-order obstructions, each
  controlling the ``deformability'' of lower-degree data.
\end{enumerate}
\end{proposition}

\begin{theorem}[Independence of cover]
\label{thm:cover-independence}
Let $\mathrm{Cov}(X)$ denote the filtered category of covers of~$X$
with refinements as morphisms.  Then the \emph{absolute \v{C}ech
cohomology}
\[
  \check{H}^n(X, \mathcal{F})
  \;=\;
  \varinjlim_{\mathcal{U} \in \mathrm{Cov}(X)}
  \check{H}^n(\mathcal{U}, \mathcal{F})
\]
is independent of the choice of cover: any two covers have a common
refinement, and the refinement maps induce isomorphisms on the
colimit.
\end{theorem}

\section{Long exact sequences}
\label{sec:les}

\begin{theorem}[Long exact sequence in cohomology]
\label{thm:les}
Let $0 \to \mathcal{F}' \to \mathcal{F} \to \mathcal{F}'' \to 0$ be a
short exact sequence of sheaves on~$(C,J)$.  There is a long exact
sequence of abelian groups:
\[
  0 \to H^0(\mathcal{F}') \to H^0(\mathcal{F}) \to H^0(\mathcal{F}'')
  \xrightarrow{\;\partial^0\;}
  H^1(\mathcal{F}') \to H^1(\mathcal{F}) \to H^1(\mathcal{F}'')
  \xrightarrow{\;\partial^1\;} H^2(\mathcal{F}') \to \cdots
\]
The \emph{connecting homomorphisms}~$\partial^n$ are constructed by
diagram chasing: given an $n$-cocycle in $\Cech^n(\mathcal{F}'')$,
lift it to $\Cech^n(\mathcal{F})$, apply the coboundary~$\delta^n$
(which lands in the image of $\mathcal{F}'$), and project to
$\Cech^{n+1}(\mathcal{F}')$.  One verifies that the result is a
cocycle and that its class is independent of the lift.
\end{theorem}

\begin{remark}
In Judgment Geometry, the long exact sequence governs how obstructions
propagate through sub-judgments.  If a judgment decomposes as
$\mathcal{F}' \hookrightarrow \mathcal{F} \twoheadrightarrow
\mathcal{F}''$ (e.g., splitting a program specification into interface
and implementation), then the connecting homomorphism~$\partial$ maps
obstructions in the quotient ($\mathcal{F}''$) to obstructions in the
sub-object ($\mathcal{F}'$).
\end{remark}

\section{Mayer--Vietoris}
\label{sec:mayer-vietoris}

\begin{theorem}[Mayer--Vietoris sequence]
\label{thm:mv}
Let $X = U \cup V$ be a decomposition of the site object into two
sub-objects.  There is a long exact sequence:
\[
  \cdots \to H^n(X,\mathcal{F})
  \xrightarrow{\;\rho\;}
  H^n(U,\mathcal{F}) \oplus H^n(V,\mathcal{F})
  \xrightarrow{\;\delta_{\mathrm{MV}}\;}
  H^n(U \cap V,\mathcal{F})
  \xrightarrow{\;\partial\;}
  H^{n+1}(X,\mathcal{F}) \to \cdots
\]
where $\rho$ is restriction and $\delta_{\mathrm{MV}}$ is the
difference of restrictions.
\end{theorem}

\begin{example}[Modular decomposition]
\label{ex:modular-mv}
A program with two interacting modules $M_1$ and $M_2$ has
$X = M_1 \cup M_2$ with $M_1 \cap M_2$ the shared interface.  The
Mayer--Vietoris sequence gives:
\[
  H^0(M_1) \oplus H^0(M_2) \to H^0(M_1 \cap M_2)
  \xrightarrow{\;\partial\;} H^1(X).
\]
Thus $H^1(X) \neq 0$---a global bug---arises precisely when the
interface specifications are individually satisfied but mutually
inconsistent.
\end{example}

\section{Spectral sequences}
\label{sec:spectral}

\begin{definition}[Spectral sequence]
\label{def:spectral}
A \emph{cohomological spectral sequence} in an abelian category is a
collection $\{E_r^{p,q},\; d_r^{p,q}\}_{r \geq 0}$ where:
\begin{itemize}[nosep]
\item Each $E_r^{p,q}$ is an object (the $(p,q)$-entry on
  page~$r$).
\item Each $d_r^{p,q} : E_r^{p,q} \to E_r^{p+r,\, q-r+1}$ is a
  differential with $d_r \circ d_r = 0$.
\item $E_{r+1}^{p,q} = H^{p,q}(E_r, d_r) =
  \ker\, d_r^{p,q} / \mathrm{im}\, d_r^{p-r,\, q+r-1}$.
\end{itemize}
The spectral sequence \emph{converges} to a graded object
$H^n = \bigoplus_{p+q=n} E_\infty^{p,q}$ if it degenerates at a
finite page: $E_r = E_\infty$ for some~$r$.
\end{definition}

\begin{theorem}[\v{C}ech-to-derived functor spectral sequence]
\label{thm:cech-derived}
For a cover $\mathcal{U}$ of~$X$ and a sheaf $\mathcal{F}$, there is a
spectral sequence with
\[
  E_2^{p,q}
  \;=\;
  \check{H}^p\!\left(\mathcal{U},\;
  \underline{H}^q(\mathcal{F})\right)
  \;\Longrightarrow\;
  H^{p+q}(X, \mathcal{F}),
\]
where $\underline{H}^q(\mathcal{F})$ is the presheaf
$U \mapsto H^q(U, \mathcal{F})$ (derived functor cohomology on open
sub-objects).  If $\mathcal{U}$ is a \emph{Leray cover}---meaning
$H^q(U_{i_0} \cap \cdots \cap U_{i_p}, \mathcal{F}) = 0$ for
$q > 0$---the spectral sequence degenerates at $E_2$ and
\[
  \check{H}^n(\mathcal{U}, \mathcal{F})
  \;\cong\; H^n(X, \mathcal{F}).
\]
\end{theorem}

\begin{remark}
In Judgment Geometry, convergence of the spectral sequence means that
\v{C}ech obstructions (computable from a cover) faithfully capture all
derived obstructions.  The Leray condition identifies ``good''
covers---those for which the combinatorial \v{C}ech computation suffices.
\end{remark}

\section{Obstruction theory}
\label{sec:obstruction}

\begin{construction}[Obstruction tower]
\label{con:obstruction-tower}
Let $X$ be built by attaching cells: $X^{(0)} \subseteq X^{(1)}
\subseteq \cdots \subseteq X$.  Given a sheaf~$\mathcal{F}$ and a
section $s_n \in \Gamma(X^{(n)}, \mathcal{F})$ over the $n$-skeleton,
the obstruction to extending $s_n$ to a section over $X^{(n+1)}$ is a
class
\[
  o_{n+1}(s_n) \;\in\; H^{n+1}(X^{(n+1)},\; X^{(n)};\; \mathcal{F}).
\]
If $o_{n+1}(s_n) = 0$, the section extends; otherwise, $s_n$ must be
modified on the $n$-skeleton (if possible) to kill the obstruction.
\end{construction}

\begin{proposition}[Systematic obstruction theory]
\label{prop:systematic-obstruction}
A global section $s \in \Gamma(X, \mathcal{F})$ exists if and only if
all obstructions $o_1, o_2, o_3, \ldots$ can be successively killed.
More precisely, an inductive construction produces:
\begin{enumerate}[label=(\alph*),nosep]
\item $s_0 \in \Gamma(X^{(0)}, \mathcal{F})$ exists freely
  (sections on a discrete set).
\item $s_n$ extends to $s_{n+1}$ iff $o_{n+1}(s_n) = 0$.
\item If $o_{n+1} \neq 0$ but represents the zero class in the
  colimit $\check{H}^{n+1}(X, \mathcal{F})$, then $s_n$ can be
  refined to kill~$o_{n+1}$.
\end{enumerate}
\end{proposition}

\section{The derived category}
\label{sec:derived}

\begin{definition}[Derived category]
\label{def:derived}
The \emph{bounded derived category} $D^b(\Sh(C,J))$ is obtained from
the category of bounded cochain complexes of sheaves by formally
inverting quasi-isomorphisms (maps inducing isomorphisms on all
cohomology sheaves).  It is a \emph{triangulated category} with
shift functor $[1]$ and distinguished triangles.
\end{definition}

\begin{remark}
The \v{C}ech complex $\Cech^\bullet(\mathcal{U}, \mathcal{F})$ is an
object of~$D^b$.  Its cohomology is $R\Gamma(X, \mathcal{F})$---the
derived global sections.  The advantage of working in $D^b$ is that all
spectral sequence arguments become statements about distinguished
triangles, and the derived Hom
$\mathrm{RHom}(\mathcal{F}, \mathcal{G})$ encodes extension classes
$\mathrm{Ext}^n(\mathcal{F}, \mathcal{G})$ that measure higher
obstructions to morphisms between sheaves.
\end{remark}

\section{Characteristic classes}
\label{sec:char-classes}

In classical topology, characteristic classes attach cohomological
invariants to fibre bundles, measuring how ``twisted'' the bundle is.
In Judgment Geometry, the role of the fibre bundle is played by
the sheaf~$\mathcal{F}$ over the site, and the characteristic classes
measure the complexity of the judgment's non-trivial topology.

\begin{definition}[Characteristic class in JuGeo]
\label{def:char-class}
Let $\mathcal{F}$ be a locally free sheaf of rank~$r$ on the site
$(C,J)$.  The \emph{$k$-th characteristic class}
$c_k(\mathcal{F}) \in H^{2k}(X, \mathbb{Z})$ (in the classical case)
or $c_k(\mathcal{F}) \in \check{H}^k(X, \mathcal{F}^\times)$ (in
the general sheaf-theoretic case) measures the degree-$k$ obstruction
to trivialising~$\mathcal{F}$.
\begin{itemize}[nosep]
\item $c_0 = 1$: the rank is locally constant.
\item $c_1$: the first obstruction to choosing a global frame---in
  JuGeo, whether local evidence can be coherently oriented.
\item Higher $c_k$: successive obstructions to higher-order
  trivialisations.
\end{itemize}
\end{definition}

\section{The complete definition}

\begin{definition}[Cohomological Spectrum $\CSpec$]
\label{def:cspec}
A \emph{Cohomological Spectrum} is a tuple
\[
  \CSpec \;=\;
  \bigl(\mathcal{U},\; \Cech^\bullet,\;
  (\check{H}^n)_{n \geq 0},\; (\delta^n),\;
  \partial,\; \mathrm{LES},\; \mathrm{MV},\;
  E_r^{p,q},\; \mathrm{char},\; D^b\bigr)
\]
where:
\begin{enumerate}[label=(\roman*),nosep]
\item $\mathcal{U}$ is a cover of the site object
  (Definition~\ref{def:cover});
\item $\Cech^\bullet = \Cech^\bullet(\mathcal{U}, \mathcal{F})$ is the
  \v{C}ech cochain complex (Definition~\ref{def:cech});
\item $(\check{H}^n)_{n \geq 0}$ are the \v{C}ech cohomology groups
  (Definition~\ref{def:cech-cohom});
\item $(\delta^n)$ are the coboundary maps
  (Definition~\ref{def:coboundary-p2});
\item $\partial$ denotes the connecting homomorphisms of the long exact
  sequence (Theorem~\ref{thm:les});
\item $\mathrm{LES}$ is the long exact sequence datum
  (Theorem~\ref{thm:les});
\item $\mathrm{MV}$ is the Mayer--Vietoris sequence datum
  (Theorem~\ref{thm:mv});
\item $E_r^{p,q}$ is the \v{C}ech-to-derived spectral sequence
  (Theorem~\ref{thm:cech-derived});
\item $\mathrm{char}$ denotes the characteristic classes
  (Definition~\ref{def:char-class});
\item $D^b = D^b(\Sh(C,J))$ is the bounded derived category
  (Definition~\ref{def:derived}).
\end{enumerate}
\end{definition}

\begin{definition}[Morphism of Cohomological Spectra]
\label{def:cspec-morph}
A \emph{morphism} $f : \CSpec_1 \to \CSpec_2$ consists of:
\begin{enumerate}[label=(\alph*),nosep]
\item a refinement of covers $\lambda : \mathcal{U}_2 \to \mathcal{U}_1$;
\item induced cochain maps
  $\lambda^* : \Cech^\bullet(\mathcal{U}_1, \mathcal{F}_1) \to
  \Cech^\bullet(\mathcal{U}_2, \mathcal{F}_2)$;
\item compatibility with all connecting homomorphisms:
  $\lambda^* \circ \partial_1 = \partial_2 \circ \lambda^*$.
\end{enumerate}
\end{definition}

We write $\mathbf{CSpec}$ for the category of Cohomological Spectra.

\section{Examples across domains}
\label{sec:cspec-examples}

\begin{example}[Program bug as $H^1$ class]
\label{ex:program-bug}
A function \texttt{merge\_sort} with two branches---one for lists of
length $\leq 1$, one for longer lists---has cover
$\mathcal{U} = \{U_{\mathrm{base}},\, U_{\mathrm{rec}}\}$.
A $0$-cochain assigns a sorting invariant to each branch; a $1$-cocycle
on the overlap $U_{\mathrm{base}} \cap U_{\mathrm{rec}}$ records
whether the invariants match at the boundary.
If $H^1 \neq 0$, the local invariants are inconsistent: there is a bug
at the case split.
\end{example}

\begin{example}[NL ambiguity as $H^1$ class]
\label{ex:nl-ambiguity}
In the sentence ``Every farmer who owns a donkey beats it,'' the
discourse site has coordinates for ``every farmer,'' ``a donkey,'' and
``it.''  The presheaf assigns possible referents at each coordinate.
The $1$-cocycle measures whether the pronoun ``it'' can be coherently
resolved across all accessible antecedents.  When
$H^1 \neq 0$, the pronoun is genuinely ambiguous---no single
assignment of referents is globally consistent.
\end{example}

\begin{example}[Poetic tension as $H^1$ class]
\label{ex:poetic-tension}
In a sonnet with an ABAB rhyme scheme, the cover consists of
quatrains.  A $0$-cochain assigns a consistent rhyme pattern to each
quatrain; a $1$-cocycle on the overlap between quatrains~1 and~2
measures whether the shared rhyme (the B-rhyme of quatrain~1 and the
A-rhyme of quatrain~2) is phonologically consistent.  When $H^1 \neq
0$, there is metrical or phonological tension between quatrains---a
feature exploited deliberately in English sonnets.
\end{example}

\begin{example}[Proof gap as $H^1$ class]
\label{ex:proof-gap}
A mathematical proof split into lemmas $L_1, \ldots, L_k$ has a cover
$\mathcal{U} = \{L_i\}$.  Each lemma's proof is a $0$-cochain; the
$1$-cocycle on the overlap $L_i \cap L_{i+1}$ records whether the
conclusion of~$L_i$ matches the hypothesis of~$L_{i+1}$.  When
$H^1 \neq 0$, there is a \emph{proof gap}---the lemmas do not chain
together into a complete proof.
\end{example}

\chapter{The Graded Evaluation \texorpdfstring{$\GEval$}{Q}}
\label{ch:geval}

The fourth object encodes the \emph{quality} of a judgment and the
\emph{history} by which that quality was established.  At the
component level, the pair $(T,\Pi)$ records a single grade and a
provenance trail.  We now develop the full algebraic structure that
governs how grades compose, how they respond to new evidence, and how
they converge under iterated re-evaluation.

\section{The grade semiring}
\label{sec:grade-semiring}

\begin{definition}[Grade semiring]
\label{def:grade-semiring-p2}
The \emph{grade semiring} is the algebraic structure
\[
  G \;=\; ([0,1],\; \opG,\; \otG,\; \gzero,\; \gone)
\]
where $[0,1]$ is the real unit interval and the operations are:
\begin{itemize}[nosep]
\item $\opG$ (join): $a \opG b = \max(a,b)$,
\item $\otG$ (tensor): $a \otG b = a \cdot b$ (multiplication),
\item $\gzero = 0$ (additive identity / worst grade),
\item $\gone = 1$ (multiplicative identity / perfect grade).
\end{itemize}
\end{definition}

\begin{proposition}[Semiring axioms]
\label{prop:semiring}
$(G, \opG, \otG, \gzero, \gone)$ is a commutative semiring:
\begin{enumerate}[label=(\alph*),nosep]
\item $\opG$ is associative, commutative, and $\gzero$ is its identity:
  $a \opG \gzero = a$.
\item $\otG$ is associative, commutative, and $\gone$ is its identity:
  $a \otG \gone = a$.
\item $\otG$ distributes over $\opG$:
  $a \otG (b \opG c) = (a \otG b) \opG (a \otG c)$.
\item $\gzero$ annihilates: $a \otG \gzero = \gzero$.
\end{enumerate}
Moreover, $\opG$ is \emph{idempotent}: $a \opG a = a$.
\end{proposition}

\begin{proof}
All properties follow from the standard properties of $\max$ and
$\cdot$ on $[0,1]$.  Idempotency of~$\opG$ is immediate:
$\max(a,a) = a$.  Distributivity: $a \cdot \max(b,c) = \max(a \cdot b,
a \cdot c)$ since multiplication by $a \geq 0$ preserves order.
\end{proof}

\section{Lattice and quantale structure}
\label{sec:quantale}

\begin{definition}[Partial order]
\label{def:partial-order}
Define $a \leq b$ iff $a \opG b = b$, i.e., $\max(a,b) = b$, i.e.,
$a \leq b$ in the usual order on $[0,1]$.
\end{definition}

\begin{theorem}[Complete lattice]
\label{thm:lattice}
$(G, \leq)$ is a complete lattice.  For any subset
$S \subseteq [0,1]$:
\[
  \bigvee S = \sup S, \qquad
  \bigwedge S = \inf S.
\]
In particular, $\opG = \bigvee$ (binary join), $\gzero = \bot$, and
$\gone = \top$.
\end{theorem}

\begin{theorem}[Quantale]
\label{thm:quantale}
$(G, \otG, \gone, \leq)$ is a \emph{commutative unital quantale}:
$\otG$ distributes over arbitrary joins:
\[
  a \otG \bigvee S \;=\; \bigvee \{a \otG s : s \in S\}.
\]
This follows from the continuity of multiplication on $[0,1]$.
\end{theorem}

\section{Residuation}
\label{sec:residuation}

\begin{definition}[Residuation]
\label{def:residuation-p2}
The \emph{residual} (or \emph{internal hom}) $\multimap$ is the right
adjoint to $\otG$ in the lattice order:
\[
  a \otG c \leq b
  \quad\Longleftrightarrow\quad
  c \leq a \multimap b.
\]
Explicitly:
\[
  a \multimap b \;=\;
  \bigvee\, \{c \in [0,1] : a \otG c \leq b\}
  \;=\;
  \begin{cases}
    1 & \text{if } a \leq b, \\
    b / a & \text{if } a > b > 0, \\
    0 & \text{if } a > 0,\; b = 0.
  \end{cases}
\]
\end{definition}

\begin{theorem}[Closed monoidal structure]
\label{thm:closed-monoidal}
$(G, \otG, \multimap, \gone)$ is a symmetric closed monoidal structure
on the lattice $(G, \leq)$.  The residuation satisfies:
\begin{enumerate}[label=(\alph*),nosep]
\item $\gone \multimap a = a$ for all~$a$.
\item $a \multimap (b \multimap c) = (a \otG b) \multimap c$
  (currying).
\item $a \leq b$ implies $b \multimap c \leq a \multimap c$
  (contravariance in the first argument).
\item $a \leq b$ implies $c \multimap a \leq c \multimap b$
  (covariance in the second argument).
\end{enumerate}
\end{theorem}

\begin{remark}[Interpretation]
The residual $a \multimap b$ answers the question: ``Given that our
current evidence has quality~$a$, what is the best quality~$b$ that
can be achieved?''  If $a \leq b$ (the evidence already meets the
target), the answer is $\gone$ (perfect).  If $a > b$, the answer
degrades proportionally.
\end{remark}

\section{Grade and proof trail}
\label{sec:grade-trail}

\begin{definition}[Grade]
\label{def:grade}
The \emph{grade} of a judgment is an element $T \in G$.  It represents
the current assessed quality of the evidence supporting the judgment.
\end{definition}

\begin{definition}[Proof trail]
\label{def:trail}
The \emph{proof trail} is a finite sequence
\[
  \Pi \;=\; [(T_1, r_1),\; (T_2, r_2),\; \ldots,\; (T_k, r_k)]
\]
where each $T_i \in G$ is a grade and each $r_i$ is a \emph{reason}---
a symbolic tag identifying the source of the assessment (e.g.,
``type-checker,'' ``runtime test,'' ``SMT solver,'' ``human review'').
\end{definition}

\begin{definition}[Trail aggregation]
\label{def:aggregation}
The grade $T$ is determined from the trail~$\Pi$ by one of two
canonical aggregation modes:
\begin{itemize}[nosep]
\item \textbf{Optimistic (join):}
  $T = \bigoplus_{i=1}^{k} T_i = \max(T_1, \ldots, T_k)$ ---
  the best assessment wins.
\item \textbf{Conservative (tensor):}
  $T = \bigotimes_{i=1}^{k} T_i = T_1 \cdot T_2 \cdots T_k$ ---
  all assessments must concur; any low grade drags the aggregate down.
\end{itemize}
The choice of mode is part of the judgment's policy.
\end{definition}

\section{The evaluation functor}
\label{sec:eval-functor}

\begin{definition}[Evaluation functor]
\label{def:eval-functor}
Let $\mathbf{Proof}_\Gamma$ be the proof category
(Construction~\ref{con:proof-category}).
The \emph{evaluation functor}
$\mu : \mathbf{Proof}_\Gamma \to G$ is a function assigning a grade to
each morphism (proof term) such that:
\begin{enumerate}[label=(\alph*),nosep]
\item $\mu(\id_A) = \gone$ for every type~$A$.
\item $\mu(g \circ f) = \mu(g) \otG \mu(f)$ for composable
  morphisms~$f$ and~$g$.
\end{enumerate}
That is, $\mu$ is a \emph{lax monoidal functor} from
$(\mathbf{Proof}_\Gamma, \circ, \id)$ to $(G, \otG, \gone)$.
\end{definition}

\begin{remark}
Condition~(b) says that the grade of a composite proof is the tensor
of the grades of its parts.  Since $a \otG b \leq \min(a,b)$ on
$[0,1]$, composing proofs can only maintain or reduce quality---the
weakest link constrains the chain.
\end{remark}

\section{Stability and convergence}
\label{sec:stability-convergence}

\begin{definition}[Stability]
\label{def:stability}
The evaluation functor~$\mu$ is \emph{$L$-stable} with respect to a
metric $d$ on the presheaf space if
\[
  |T(\varphi) - T(\varphi')| \;\leq\; L \cdot d(\varphi, \varphi')
\]
for all presheaves $\varphi, \varphi'$.  Small perturbations of the
input data produce at most proportionally small changes in the grade.
\end{definition}

\begin{definition}[Convergence]
\label{def:convergence}
Define the \emph{re-evaluation sequence}
\[
  T^{(0)} = T, \qquad
  T^{(n+1)} = \mu\!\left(\mathrm{refine}(E, T^{(n)})\right),
\]
where $\mathrm{refine}(E, T^{(n)})$ re-evaluates the evidence~$E$ in
light of the grade $T^{(n)}$.  The evaluation is \emph{convergent} if
this sequence has a limit $T^* = \lim_{n \to \infty} T^{(n)}$.
\end{definition}

\begin{theorem}[Fixpoint convergence]
\label{thm:fixpoint}
If the re-evaluation map $T \mapsto \mu(\mathrm{refine}(E, T))$ is
monotone and continuous on the complete lattice $(G, \leq)$, then by
the Kleene fixed-point theorem the re-evaluation sequence converges
to the least fixpoint:
\[
  T^* \;=\; \bigvee_{n \geq 0}\, T^{(n)}.
\]
This is the \FMLM{} fixpoint property: iterated re-evaluation
stabilises.
\end{theorem}

\begin{remark}
Convergence is essential in practice: a judgment's grade should not
oscillate indefinitely under re-evaluation.  The monotonicity
condition says that higher grades cannot lead to lower re-evaluated
grades---a natural requirement when refinement improves evidence.
\end{remark}

\section{Entropy and information content}
\label{sec:entropy}

\begin{definition}[Grade entropy]
\label{def:entropy}
The \emph{entropy} of a grade $T \in (0,1)$ is
\[
  \mathcal{E}(T) \;=\;
  -T \ln T \;-\; (1-T) \ln(1-T),
\]
with $\mathcal{E}(0) = \mathcal{E}(1) = 0$ by continuity.
\end{definition}

\begin{proposition}[Entropy properties]
\label{prop:entropy}
\hfill
\begin{enumerate}[label=(\alph*),nosep]
\item $\mathcal{E}(T) \geq 0$ for all $T \in [0,1]$.
\item $\mathcal{E}$ is maximised at $T = 1/2$ (maximal uncertainty).
\item $\mathcal{E}(T) \to 0$ as $T \to 0$ or $T \to 1$ (confident
  grades carry low entropy).
\item $\mathcal{E}$ is concave.
\end{enumerate}
The entropy measures the \emph{information content} of the grade: a
grade near $0$ or $1$ is informative (high confidence), while a grade
near $1/2$ is uninformative (maximal uncertainty).
\end{proposition}

\section{Bayesian update}
\label{sec:bayesian}

\begin{definition}[Grade update]
\label{def:bayes}
When new evidence $e$ arrives with likelihood $\ell(e \mid \varphi)$
and prior grade~$T$, the \emph{updated grade} is
\[
  T' \;=\;
  \frac{T \otG \ell(e \mid \varphi)}
       {T \otG \ell(e \mid \varphi) \;+\;
        (1 - T) \otG \ell(e \mid \neg\varphi)},
\]
which is the standard Bayesian update when $\otG$ is multiplication.
The proof trail extends by one entry: $\Pi' = \Pi \cdot [(T', e)]$.
\end{definition}

\begin{proposition}
\label{prop:bayes-coherence}
Bayesian update is coherent with the quantale structure:
\begin{enumerate}[label=(\alph*),nosep]
\item If $\ell(e \mid \varphi) = \gone$, then $T' \geq T$: perfectly
  supporting evidence can only increase the grade.
\item If $\ell(e \mid \varphi) = \gzero$, then $T' = \gzero$:
  contradicting evidence kills the grade.
\item The update is monotone in both $T$ and $\ell$.
\end{enumerate}
\end{proposition}

\section{The trust algebra}
\label{sec:trust-algebra}

\begin{definition}[Trust composition]
\label{def:trust-composition-p2}
For a chain of judgments $J_1 \circ J_2$ (judgment $J_2$ depends on the
conclusion of~$J_1$), the composite grade is
\[
  T(J_1 \circ J_2) \;=\; T(J_1) \otG T(J_2).
\]
Since $\otG$ is multiplication on $[0,1]$, trust degrades along chains.
\end{definition}

\begin{proposition}[Sub-transitivity]
\label{prop:sub-transitive}
Trust is \emph{sub-transitive} in general:
\[
  T(A \to B) \otG T(B \to C) \;\leq\; T(A \to C).
\]
Equality holds when the intermediate judgment~$B$ introduces no
additional uncertainty (i.e., $T(B \to C)$ already accounts for the
quality of~$B$).  In general, the direct path $A \to C$ may have
higher trust than the indirect path through~$B$, because it avoids
compounding errors.
\end{proposition}

\begin{definition}[Grade polynomials]
\label{def:grade-poly}
A \emph{grade polynomial} is an algebraic expression built from grades
using $\opG$ and $\otG$:
\[
  T \;=\; (T_1 \otG T_2) \opG (T_3 \otG T_4).
\]
Such expressions arise when a judgment has multiple independent proof
paths (joined with~$\opG$), each of which is a chain of sub-judgments
(tensored with~$\otG$).  Evaluation of grade polynomials in the
quantale gives the aggregate grade.
\end{definition}

\section{The Eilenberg--Moore category}
\label{sec:eilenberg-moore}

\begin{definition}[$G$-modules]
\label{def:g-module}
A \emph{$G$-module} is a set $M$ equipped with a $G$-action
$\bullet : G \times M \to M$ satisfying:
\begin{enumerate}[label=(\alph*),nosep]
\item $\gone \bullet m = m$ for all $m \in M$.
\item $(a \otG b) \bullet m = a \bullet (b \bullet m)$.
\item $a \bullet (m \opG_M n) = (a \bullet m) \opG_M (a \bullet n)$,
  where $\opG_M$ is a join operation on~$M$.
\end{enumerate}
The category $G\text{-}\mathbf{Mod}$ of $G$-modules and equivariant
maps is the natural home for graded data: each element of~$M$
represents a datum, and the $G$-action scales its quality.
\end{definition}

\begin{remark}
Graded judgments form a $G$-module: the underlying set is the set of
judgments, and the action $T \bullet \JJ$ re-grades $\JJ$ by
tensoring with~$T$.  Morphisms of $G$-modules are grade-compatible
judgment transformations.
\end{remark}

\section{The complete definition}

\begin{definition}[Graded Evaluation $\GEval$]
\label{def:geval}
A \emph{Graded Evaluation} is a tuple
\[
  \GEval \;=\;
  \bigl(G,\; \opG,\; \otG,\; \gzero,\; \gone,\;
  {\leq},\; T,\; \Pi,\; {\multimap},\;
  \mu,\; \sigma,\; \kappa,\; \mathcal{E},\;
  \mathrm{Bayes}\bigr)
\]
where:
\begin{enumerate}[label=(\roman*),nosep]
\item $(G, \opG, \otG, \gzero, \gone)$ is the grade semiring
  (Definition~\ref{def:grade-semiring-p2});
\item $\leq$ is the induced partial order
  (Definition~\ref{def:partial-order});
\item $T \in G$ is the current grade (Definition~\ref{def:grade});
\item $\Pi$ is the proof trail (Definition~\ref{def:trail});
\item $\multimap$ is the residuation
  (Definition~\ref{def:residuation-p2});
\item $\mu$ is the evaluation functor
  (Definition~\ref{def:eval-functor});
\item $\sigma$ is the stability constant
  (Definition~\ref{def:stability});
\item $\kappa$ is the convergence datum
  (Definition~\ref{def:convergence});
\item $\mathcal{E}$ is the entropy function
  (Definition~\ref{def:entropy});
\item $\mathrm{Bayes}$ is the Bayesian update rule
  (Definition~\ref{def:bayes}).
\end{enumerate}
\end{definition}

\begin{definition}[Morphism of Graded Evaluations]
\label{def:geval-morph}
A \emph{morphism} $f : \GEval_1 \to \GEval_2$ is a semiring
homomorphism $f_G : G_1 \to G_2$ together with a map of trails
$f_\Pi : \Pi_1 \to \Pi_2$ such that:
\begin{enumerate}[label=(\alph*),nosep]
\item $f_G(T_1) = T_2$;
\item $f_G$ preserves the residuation:
  $f_G(a \multimap b) = f_G(a) \multimap f_G(b)$;
\item $f_\Pi$ is order-preserving and respects the evaluation functor:
  $\mu_2 \circ f_\Pi = f_G \circ \mu_1$.
\end{enumerate}
\end{definition}

We write $\mathbf{GEval}$ for the category of Graded Evaluations.

\section{Examples}

\begin{example}[Test coverage]
\label{ex:test-coverage}
A program's test suite produces a Graded Evaluation.  Each test~$t_i$
yields a grade $T_i \in \{0, 1\}$ (pass/fail); the aggregate is
$T = \bigotimes_i T_i$ under conservative aggregation (all tests must
pass) or $T = |\{i : T_i = 1\}| / k$ (coverage fraction) under a
softer policy.  The trail
$\Pi = [(T_1, t_1), \ldots, (T_k, t_k)]$ records which tests
contributed.  The entropy $\mathcal{E}(T)$ is high for $T \approx
0.5$ (half the tests fail), indicating maximal uncertainty about
correctness.
\end{example}

\begin{example}[Semantic adequacy]
\label{ex:semantic-adequacy}
A semantic parse of a sentence receives a grade $T$ from a model that
estimates how well the $\lambda$-term captures the intended meaning.
The trail records: the parser's confidence ($T_1$), human evaluation
($T_2$), and entailment test ($T_3$).  Bayesian update incorporates
each source sequentially, and the entropy decreases as evidence
accumulates.
\end{example}

\begin{example}[Poetic quality]
\label{ex:poetic-quality}
A poem's quality is graded by multiple criteria: meter regularity
($T_1$), rhyme consistency ($T_2$), semantic coherence ($T_3$), and
aesthetic judgement ($T_4$).  The grade polynomial
$T = (T_1 \otG T_2) \opG (T_3 \otG T_4)$ says: the overall quality
is the best of two independent assessments---formal (meter $\otG$
rhyme) and semantic (coherence $\otG$ aesthetics).
\end{example}

\begin{example}[Proof confidence]
\label{ex:proof-confidence}
A formally verified theorem has $T = \gone$ by a single trail entry
$(1, \text{proof-checker})$.  A theorem verified by an SMT solver with
a timeout has $T = 0.95$ with reason ``solver-with-timeout.''  If later
a formal proof is produced, the trail grows by $(\gone,
\text{proof-checker})$ and the optimistic aggregate becomes $T = \gone$.
\end{example}

\chapter{The Judgment Fibration}
\label{ch:fibration}

The four objects $\STop$, $\PStr$, $\CSpec$, $\GEval$ do not stand
independently: they are bound together by compatibility conditions that
express how the semantic, proof-theoretic, cohomological, and
evaluative aspects of a judgment constrain one another.  This chapter
formulates these constraints as a \emph{fibration} over the category
of Semantic Topoi, making precise the idea that the other three objects
are ``fibred'' over the site.

\section{Compatibility conditions}
\label{sec:compatibility}

\begin{definition}[Compatible quadruple]
\label{def:compatible}
A \emph{compatible quadruple} $(\STop, \PStr, \CSpec, \GEval)$
consists of one object from each of the four categories, subject to the
following conditions:

\medskip
\noindent\textbf{(C1) Topos determines context.}
The context $\Gamma$ of the Proof Structure is drawn from the internal
language of the Semantic Topos:
\begin{itemize}[nosep]
\item Each variable declaration $x_i : A_i$ in~$\Gamma$ corresponds to
  a generalised element $x_i : 1 \to \varphi_i$ in the
  Mitchell--B\'enabou language of~$\STop$.
\item The type $A$ is an object of the internal type theory of the
  topos.
\item Well-formedness of~$\Gamma$ is equivalent to well-definedness
  of the corresponding string of internal terms.
\end{itemize}

\medskip
\noindent\textbf{(C2) Proof determines cohomology.}
The cover $\mathcal{U}$ of the Cohomological Spectrum arises from
the proof search for~$E : A$:
\begin{itemize}[nosep]
\item The cover $\mathcal{U} = \{U_i\}$ consists of the sub-goals into
  which the proof obligation $\Gamma \vdash ? : A$ is decomposed.
\item A $1$-cocycle in $\check{H}^1(\mathcal{U}, \mathcal{F})$ records
  a failure to glue sub-proofs: the boundary
  $(\delta^0 E)_{ij} = E_i|_{U_{ij}} - E_j|_{U_{ij}}$ does not
  vanish.
\item In particular: $\check{H}^1 = 0$ if and only if the sub-proofs
  glue into a global proof $E$, recovering the fundamental theorem of
  Judgment Geometry.
\end{itemize}

\medskip
\noindent\textbf{(C3) Cohomology constrains grade.}
The grade~$T$ of the Graded Evaluation is bounded by the cohomological
data:
\begin{itemize}[nosep]
\item If $\check{H}^1(\mathcal{U}, \mathcal{F}) = 0$, then the grade
  can be high: $T$ depends only on the quality of~$E$.
\item If $\check{H}^1 \neq 0$, then $T$ is bounded:
  \[
    T \;\leq\; \gone \;-\; \|[\check{H}^1]\|
  \]
  for some norm $\|\cdot\|$ on the cohomology group.  Non-trivial
  obstructions cap the achievable quality.
\item Higher cohomology provides finer bounds:
  $\check{H}^2 \neq 0$ implies $T$ is further reduced, reflecting
  the impossibility of even modifying the attempted gluing.
\end{itemize}

\medskip
\noindent\textbf{(C4) Grade feeds back to topos.}
The Graded Evaluation informs the construction of future Semantic
Topoi:
\begin{itemize}[nosep]
\item Low-grade regions of the site may be refined (adding new
  coordinates) or coarsened (merging coordinates) to improve
  future judgments.
\item The Lawvere--Tierney topology $\tau$ may be adjusted:
  replacing~$\tau$ by a finer topology (more covering sieves)
  demands stronger evidence, while a coarser topology is more
  permissive.
\item This feedback loop gives Judgment Geometry its
  \emph{adaptive} character.
\end{itemize}
\end{definition}

\section{The fibration}
\label{sec:fibration-def}

\begin{construction}[Total category of judgments]
\label{con:total-category}
Define the category $\JJ$ of \emph{judgments} as follows:
\begin{itemize}[nosep]
\item \textbf{Objects:} compatible quadruples
  $(\STop, \PStr, \CSpec, \GEval)$ as in
  Definition~\ref{def:compatible}.
\item \textbf{Morphisms:} a morphism
  $(f, \sigma, \lambda, h) :
  (\STop_1, \PStr_1, \CSpec_1, \GEval_1) \to
  (\STop_2, \PStr_2, \CSpec_2, \GEval_2)$
  consists of:
  \begin{enumerate}[label=(\alph*),nosep]
  \item a morphism $f : \STop_1 \to \STop_2$ of Semantic Topoi
    (Definition~\ref{def:stop-morph});
  \item a morphism $\sigma : \PStr_1 \to \PStr_2$ of Proof Structures
    (Definition~\ref{def:pstr-morph}), compatible with~$f$ via~(C1);
  \item a morphism $\lambda : \CSpec_1 \to \CSpec_2$ of Cohomological
    Spectra (Definition~\ref{def:cspec-morph}), compatible with
    $\sigma$ via~(C2);
  \item a morphism $h : \GEval_1 \to \GEval_2$ of Graded Evaluations
    (Definition~\ref{def:geval-morph}), compatible with $\lambda$
    via~(C3).
  \end{enumerate}
\item \textbf{Composition:} componentwise.
\item \textbf{Identity:} $(\id_\STop, \id_\PStr, \id_\CSpec,
  \id_\GEval)$.
\end{itemize}
\end{construction}

\begin{definition}[Judgment fibration]
\label{def:judgment-fibration}
The \emph{judgment fibration} is the functor
\[
  \pi : \JJ \;\longrightarrow\; \mathbf{STop}
\]
defined by $\pi(\STop, \PStr, \CSpec, \GEval) = \STop$ on objects and
$\pi(f, \sigma, \lambda, h) = f$ on morphisms.  This functor
``projects'' a judgment onto its underlying Semantic Topos, forgetting
the proof, cohomological, and evaluative data.
\end{definition}

\begin{theorem}[Fibration structure]
\label{thm:fibration}
The functor $\pi : \JJ \to \mathbf{STop}$ is a Grothendieck fibration:
for every morphism $f : \STop_1 \to \STop_2$ in~$\mathbf{STop}$ and
every object $(\STop_2, \PStr_2, \CSpec_2, \GEval_2)$ in the fibre
over~$\STop_2$, there exists a \emph{Cartesian lift}---a morphism
in~$\JJ$ over~$f$ that is universal among all lifts.
\end{theorem}

\begin{proof}[Proof sketch]
The Cartesian lift is constructed by pulling back the proof, cohomological,
and evaluative data along the geometric morphism~$f$.
Specifically:
\begin{enumerate}[label=(\roman*),nosep]
\item The context $\Gamma_1$ is obtained by applying $f^*$ to the
  types in~$\Gamma_2$.
\item The cover $\mathcal{U}_1$ is the pullback cover
  $f^*(\mathcal{U}_2)$.
\item The grade $T_1 = T_2$ is unchanged (grades are scalar).
\end{enumerate}
The universality follows from the universality of pullback in each
component category.
\end{proof}

\section{Fibres and sections}
\label{sec:fibres-sections}

\begin{definition}[Fibre]
The \emph{fibre} over a Semantic Topos $\STop$ is the subcategory
\[
  \pi^{-1}(\STop)
  \;=\;
  \{(\PStr, \CSpec, \GEval) :
  (\STop, \PStr, \CSpec, \GEval) \in \ob(\JJ)\}
\]
with morphisms $(\id_\STop, \sigma, \lambda, h)$.
\end{definition}

The fibre $\pi^{-1}(\STop)$ is the space of all possible judgments
\emph{over a fixed site and presheaf}.  It encodes the freedom in
choosing evidence, decomposing into sub-goals, and assigning grades.

\begin{definition}[Section]
\label{def:section}
A \emph{section} of the fibration~$\pi$ is a functor
$s : \mathbf{STop} \to \JJ$ with $\pi \circ s = \id_{\mathbf{STop}}$.
That is, $s$ assigns to each Semantic Topos a compatible judgment
over it, coherently across all morphisms of topoi.
\end{definition}

\begin{remark}
A section is a \emph{global judgment policy}: for every site and
presheaf, it specifies how to construct evidence, decompose into
sub-goals, and assign grades, in a way that is compatible with
all changes of site.
\end{remark}

\section{Harmony as section}
\label{sec:harmony-section}

\begin{definition}[Harmonious judgment]
\label{def:harmonious}
A judgment $(\STop, \PStr, \CSpec, \GEval)$ is \emph{harmonious} if:
\begin{enumerate}[label=(\alph*),nosep]
\item It is compatible (Definition~\ref{def:compatible}).
\item $\check{H}^1(\mathcal{U}, \mathcal{F}) = 0$: no gluing
  obstructions.
\item $T \geq T_{\mathrm{threshold}}$ for a chosen threshold
  $T_{\mathrm{threshold}} \in G$.
\end{enumerate}
A section $s : \mathbf{STop} \to \JJ$ is \emph{harmonious} if $s(\STop)$
is harmonious for every~$\STop$.
\end{definition}

\begin{proposition}
\label{prop:harmony-section}
A program~$P$ satisfies its specification~$\varphi$ if and only if there
exists a harmonious section $s$ of the judgment fibration with
$s(\STop_P) = (\STop_P, \PStr, \CSpec, \GEval)$ where $\PStr$ contains
a valid proof $E : A$ and $T = \gone$.
\end{proposition}

This recovers the fundamental theorem of Judgment Geometry: correctness
is the existence of a harmonious section with maximal grade.

\section{The bottleneck principle}
\label{sec:bottleneck-principle}

\begin{theorem}[Bottleneck principle]
\label{thm:bottleneck-p2}
For a section $s$ over a cover $\{U_i\}$ of the site, the global grade
is bounded by:
\[
  T_{\mathrm{global}}
  \;=\;
  \bigotimes_{i} T_{s(U_i)}
  \;\leq\;
  \min_i \, T_{s(U_i)}.
\]
That is, the global grade is dominated by the \emph{worst local grade}.
\end{theorem}

\begin{proof}
Since $\otG$ is multiplication on $[0,1]$ and all factors are in
$[0,1]$, we have $\prod_i T_i \leq \min_i T_i$.  The product
$\bigotimes_i T_{s(U_i)}$ computes the conservative aggregate
(Definition~\ref{def:aggregation}), which is the relevant measure for
global validity.
\end{proof}

\begin{corollary}
\label{cor:bottleneck}
A single stratum with $T_{s(U_j)} = \gzero$ forces
$T_{\mathrm{global}} = \gzero$, regardless of the grades on all other
strata.  In practice: a single failing module blocks the entire
program's verification.
\end{corollary}

\section{Descent as optimisation}
\label{sec:descent-optimisation}

The judgment fibration frames the central problem of Judgment Geometry
as an \emph{optimisation problem over sections}:

\begin{definition}[Judgment optimisation]
\label{def:judgment-opt}
Given a Semantic Topos $\STop$, the \emph{judgment optimisation
problem} is to find a section $s^*$ of $\pi$ that maximises the
global grade subject to the cohomological constraint:
\[
  s^* \;=\;
  \arg\max_{s \in \Gamma(\pi)}\;
  T_{\mathrm{global}}(s)
  \qquad\text{subject to}\qquad
  \check{H}^1(\mathcal{U}_{s(\STop)},\; \mathcal{F}) = 0.
\]
\end{definition}

\begin{remark}
The constraint $\check{H}^1 = 0$ is a necessary condition for
correctness; the optimisation over~$T$ selects the highest-quality
proof among all valid ones.  When no section satisfies the constraint,
the optimum is $T^* = \gzero$, signalling that the judgment is
\emph{not} harmonious.
\end{remark}

\begin{remark}[Descent = gluing sections]
The name ``descent'' is apt: in algebraic geometry, descent theory asks
when a local datum (defined on a cover) can be glued into a global
datum.  In Judgment Geometry, a ``local judgment'' is a section defined
on each $U_i$; \emph{descent} is the process of gluing these local
judgments into a global section.  The obstruction to descent is precisely
$\check{H}^1$, and successful descent is equivalent to finding an
optimal global section.
\end{remark}

\section{Summary}

The judgment fibration $\pi : \JJ \to \mathbf{STop}$ provides a
unified geometric framework for the four objects:

\begin{center}
\renewcommand{\arraystretch}{1.3}
\begin{tabular}{@{}ll@{}}
\textbf{Concept} & \textbf{Fibration interpretation} \\
\hline
Judgment & Object of the total category $\JJ$ \\
Semantic context & Object of the base $\mathbf{STop}$ \\
Evidence + obstructions + grade & Object in the fibre $\pi^{-1}(\STop)$ \\
Correctness & Existence of a harmonious section \\
Verification & Search for a section with $T \geq T_{\mathrm{threshold}}$ \\
Bug & Obstruction: $\check{H}^1 \neq 0$ or $T < T_{\mathrm{threshold}}$ \\
Refactoring & Geometric morphism in $\mathbf{STop}$ \\
Quality & Global grade $T_{\mathrm{global}} = \bigotimes_i T_i$
\end{tabular}
\end{center}

\chapter{The Component-Wise View}
\label{ch:equivalence}

This chapter makes precise the relationship between the structured
4-tuple formulation $(\STop, \PStr, \CSpec, \GEval)$ and its
component-wise flattening $(c, \varphi, A, E, O, B, T, \Pi)$.
We construct two functors---a
flattening functor that extracts the eight components from a compatible
quadruple, and a structuring functor that assembles components into their
canonical structured objects---and prove that they form an adjoint
equivalence of categories.  The upshot: the component-wise view is a
faithful, lossless projection of the structured view.

\section{The categories}

\begin{definition}[The category $\mathbf{Judg}_8$]
\label{def:judg8}
The category $\mathbf{Judg}_8$ of \emph{component-wise judgments} has:
\begin{itemize}[nosep]
\item \textbf{Objects:} component tuples
  $\mathcal{J} = (c, \varphi, A, E, O, B, T, \Pi)$ satisfying the
  coherence axioms of Judgment Geometry.
\item \textbf{Morphisms:} a morphism
  $\mathcal{J}_1 \to \mathcal{J}_2$ is a tuple of component maps
  $(f_c, f_\varphi, f_A, f_E, f_O, f_B, f_T, f_\Pi)$ preserving
  all compatibility conditions: $f_c$ maps coordinates to coordinates,
  $f_\varphi$ is a natural transformation of presheaves compatible
  with $f_c$, $f_A$ is a type map compatible with $f_E$, and so on.
\end{itemize}
\end{definition}

\begin{definition}[The category $\mathbf{Judg}_4$]
\label{def:judg4}
The category $\mathbf{Judg}_4$ is the category $\JJ$ of
Construction~\ref{con:total-category}: compatible quadruples
$(\STop, \PStr, \CSpec, \GEval)$ with the morphisms defined there.
\end{definition}

\section{The flattening functor}
\label{sec:forgetful-functor}

\begin{construction}[Flattening functor $U$]
\label{con:forgetful-functor}
The \emph{flattening functor}
$U : \mathbf{Judg}_4 \to \mathbf{Judg}_8$ extracts the eight
components from their rich ambient structures:
\[
  U(\STop, \PStr, \CSpec, \GEval)
  \;=\;
  (c,\; \varphi,\; A,\; E,\; O,\; B,\; T,\; \Pi)
\]
where:
\begin{enumerate}[label=(\roman*),nosep]
\item $c$ is any chosen object of the underlying site category~$C$
  of~$\STop$ (or the terminal object of~$C$ if it exists).
\item $\varphi$ is the judgment presheaf from~$\STop$
  (Definition~\ref{def:phi}).
\item $A$ is the type from~$\PStr$ (component~(ii) of
  Definition~\ref{def:pstr}).
\item $E$ is the evidence term from~$\PStr$ (component~(iii)).
\item $O$ is the cover $\mathcal{U}$ from~$\CSpec$
  (component~(i) of Definition~\ref{def:cspec}).
\item $B$ is the first cohomology class
  $[\check{H}^1(\mathcal{U}, \mathcal{F})]$ from~$\CSpec$
  (component~(iii), degree $n=1$).
\item $T$ is the grade from~$\GEval$
  (component~(iii) of Definition~\ref{def:geval}).
\item $\Pi$ is the proof trail from~$\GEval$ (component~(iv)).
\end{enumerate}
On morphisms, $U(f, \sigma, \lambda, h)$ extracts the component maps
in the obvious way.
\end{construction}

\begin{proposition}
\label{prop:U-well-defined}
The functor~$U$ is well-defined: the component tuple
$U(\STop, \PStr, \CSpec, \GEval)$ satisfies all coherence axioms of
Judgment Geometry.
\end{proposition}

\begin{proof}
The compatibility conditions (C1)--(C4) of
Definition~\ref{def:compatible} imply:
\begin{itemize}[nosep]
\item $\varphi$ is a presheaf on a site containing~$c$: the pair
  $(c, \varphi)$ is well-formed.
\item $\Gamma \vdash E : A$: the pair $(A, E)$ is a valid typing
  judgment.
\item The cover~$O$ and class~$B$ arise from the same site and
  sheaf: $(O, B)$ is a valid cohomological datum.
\item The grade~$T$ is in the grade semiring with trail~$\Pi$:
  $(T, \Pi)$ is a valid graded evaluation.
\item The cross-component axioms (e.g., $B = 0 \Leftrightarrow$
  $E$ is globally defined) are guaranteed by~(C2).
  \qedhere
\end{itemize}
\end{proof}

\section{The structuring functor}
\label{sec:enrichment-functor}

\begin{construction}[Structuring functor $R$]
\label{con:enrichment-functor}
The \emph{structuring functor}
$R : \mathbf{Judg}_8 \to \mathbf{Judg}_4$ assembles each component
tuple into its canonical rich structure:
\[
  R(c, \varphi, A, E, O, B, T, \Pi)
  \;=\;
  (\STop_{\min},\; \PStr_{\min},\; \CSpec_{\min},\; \GEval_{\min})
\]
where each component is the \emph{minimal} (or \emph{free}) rich object
containing the given data:

\medskip
\noindent\textbf{(i) Minimal topos $\STop_{\min}$.}
Let $C_{c,\varphi}$ be the full subcategory of the site generated by
the coordinate~$c$ and all coordinates appearing in the support
of~$\varphi$.  Equip it with the induced Grothendieck topology.  Then
$\STop_{\min}$ is the sheaf topos $\Sh(C_{c,\varphi}, J)$ with
presheaf~$\varphi$, together with all the data of
Definition~\ref{def:stop} (which are determined by the site).
The Lawvere--Tierney topology is the trivial one ($j = \id_\Omega$).

\medskip
\noindent\textbf{(ii) Minimal proof structure $\PStr_{\min}$.}
The context~$\Gamma$ is the free context generated by the variables
occurring in~$E$, with types assigned by the minimal type theory in
which $E : A$ is derivable.  The derivation~$\mathcal{D}$ is the
canonical derivation of $\Gamma \vdash E : A$.  The reduction theory
is the standard $\beta\eta$-reduction.  The universe level is the
smallest $i$ such that $A : \Type_i$.  Type formers ($\Pi$, $\Sigma$,
$\mathrm{Id}$, $W$) are included precisely when required by~$E$.

\medskip
\noindent\textbf{(iii) Minimal cohomological spectrum $\CSpec_{\min}$.}
The cover is~$O$.  The \v{C}ech complex $\Cech^\bullet(O, \mathcal{F})$
is computed directly.  The cohomology groups $\check{H}^n$ are the
cohomology of this complex.  The long exact sequence, Mayer--Vietoris,
spectral sequence, characteristic classes, and derived category are all
determined by the site and sheaf.

\medskip
\noindent\textbf{(iv) Minimal graded evaluation $\GEval_{\min}$.}
The grade semiring is the standard $([0,1], \max, \cdot, 0, 1)$.
The grade is~$T$, the trail is~$\Pi$.  The residuation, evaluation
functor, stability, convergence, entropy, and Bayesian update are
determined by the semiring and the data.

\medskip
\noindent On morphisms, $R$ maps a component-wise morphism of tuples
to the induced morphism of rich structures.
\end{construction}

\begin{proposition}
\label{prop:R-well-defined}
The enrichment functor~$R$ is well-defined: the quadruple
$R(c, \varphi, A, E, O, B, T, \Pi)$ satisfies the compatibility
conditions (C1)--(C4) of Definition~\ref{def:compatible}.
\end{proposition}

\begin{proof}
The coherence axioms of the component tuple ensure that the components
are mutually
consistent.  The minimal constructions preserve this consistency:
(C1)~holds because the context is read off from the same site and
presheaf; (C2)~holds because the cover $O$ arose from the same proof
decomposition; (C3)~holds because the coherence axiom
$B = 0 \Leftrightarrow$ global proof exists translates directly;
(C4)~holds trivially since no feedback is active in the minimal
construction.
\end{proof}

\section{The equivalence theorem}
\label{sec:equivalence-theorem}

\begin{theorem}[Adjoint equivalence]
\label{thm:equivalence}
The functors $R : \mathbf{Judg}_8 \to \mathbf{Judg}_4$ and
$U : \mathbf{Judg}_4 \to \mathbf{Judg}_8$ form an adjoint equivalence:
\[
  R \;\dashv\; U,
  \qquad
  U \circ R \;\cong\; \id_{\mathbf{Judg}_8},
  \qquad
  R \circ U \;\cong\; \id_{\mathbf{Judg}_4}.
\]
\end{theorem}

\begin{proof}
We construct the unit and counit natural isomorphisms.

\medskip
\noindent\textbf{Unit: $\eta : \id_{\mathbf{Judg}_8}
\xrightarrow{\;\sim\;} U \circ R$.}
For a component tuple
$\mathcal{J} = (c, \varphi, A, E, O, B, T, \Pi)$:
\[
  U(R(\mathcal{J}))
  = U(\STop_{\min}, \PStr_{\min}, \CSpec_{\min}, \GEval_{\min})
  = (c', \varphi, A, E, O, B', T, \Pi).
\]
The coordinate $c' = c$ (since $c \in \ob(C_{c,\varphi})$) and the
obstruction class $B' = B$ (since the \v{C}ech cohomology of the
minimal complex with cover~$O$ yields the same $\check{H}^1$).
Thus $U \circ R(\mathcal{J}) = \mathcal{J}$, and $\eta$ is the
identity natural transformation.

\medskip
\noindent\textbf{Counit: $\varepsilon : R \circ U
\xrightarrow{\;\sim\;} \id_{\mathbf{Judg}_4}$.}
For a quadruple $\mathcal{Q} = (\STop, \PStr, \CSpec, \GEval)$:
\[
  R(U(\mathcal{Q}))
  = R(c, \varphi, A, E, O, B, T, \Pi)
  = (\STop_{\min}, \PStr_{\min}, \CSpec_{\min}, \GEval_{\min}).
\]
The minimal topos $\STop_{\min}$ is a sub-topos of~$\STop$; the
inclusion $\iota : \STop_{\min} \hookrightarrow \STop$ is a geometric
morphism.  Similarly, $\PStr_{\min}$ is a sub-structure of~$\PStr$
(the minimal type theory embeds into the ambient one), $\CSpec_{\min}$
embeds into~$\CSpec$, and $\GEval_{\min}$ embeds into~$\GEval$.

The counit $\varepsilon_\mathcal{Q} : R(U(\mathcal{Q})) \to \mathcal{Q}$
is the canonical inclusion of the minimal enrichment into the full
enrichment.  It is an isomorphism up to the choice of ambient structure:
the additional data in~$\STop$ beyond~$\STop_{\min}$ (extra objects in
the site, a non-trivial Lawvere--Tierney topology) are determined by
the site and play no role in the component-level data.  Formally,
$\varepsilon$ is a natural isomorphism because any two enrichments of
the same component data are canonically isomorphic as objects of
$\mathbf{Judg}_4$ (the enrichment is determined up to unique
isomorphism by the data it enriches).

\medskip
\noindent\textbf{Triangle identities.}
We verify $U \varepsilon \circ \eta U = \id_U$ and
$\varepsilon R \circ R \eta = \id_R$.  Both follow immediately from
$\eta = \id$ and $\varepsilon$ being the canonical inclusion.
\end{proof}

\section{Transfer of theorems}

\begin{corollary}[Theorem transfer]
\label{cor:theorem-transfer}
Every theorem about component-wise judgments transfers to structured
judgments and vice versa.  Specifically, for any property~$P$ definable
in the language of either category:
\begin{enumerate}[label=(\alph*),nosep]
\item $\mathcal{J} \in \mathbf{Judg}_8$ satisfies~$P$ if and only if
  $R(\mathcal{J}) \in \mathbf{Judg}_4$ satisfies the transported
  property $R^*(P)$.
\item $\mathcal{Q} \in \mathbf{Judg}_4$ satisfies~$Q$ if and only if
  $U(\mathcal{Q}) \in \mathbf{Judg}_8$ satisfies $U^*(Q)$.
\end{enumerate}
\end{corollary}

\begin{proof}
An adjoint equivalence is, in particular, an equivalence of categories.
Equivalent categories satisfy the same categorical properties: any
diagram that commutes in one commutes in the other, any limit or
colimit that exists in one has a counterpart in the other, and any
functor defined on one extends to the other.
\end{proof}

\begin{corollary}[Fundamental theorem, 4-tuple version]
\label{cor:fundamental}
A program~$P$ satisfies its specification~$\varphi$ if and only if
the Cohomological Spectrum of the associated judgment satisfies
$\check{H}^1(\mathcal{U}, \mathcal{F}_\varphi) = 0$.
\end{corollary}

\begin{proof}
Apply the theorem transfer to the fundamental theorem of
Judgment Geometry.
\end{proof}

\section{Strict gain in expressiveness}
\label{sec:strict-gain}

\begin{remark}[The structured view sees more]
\label{rem:strict-gain}
While the two formulations are categorically equivalent (every theorem
transfers), the structured 4-tuple is \emph{strictly more expressive in
practice} because it makes visible:

\begin{enumerate}[label=(\roman*),nosep]
\item \textbf{Morphisms.}  The component-wise view has no intrinsic
  notion of ``morphism of judgments''; the 4-tuple inherits geometric
  morphisms, context morphisms, chain maps, and grade homomorphisms
  from its ambient categories.  These morphisms enable the study of
  \emph{judgment transformations}: refactoring, generalisation,
  specialisation, composition.

\item \textbf{Internal logic.}  The Semantic Topos carries the
  Mitchell--B\'enabou language; the Proof Structure carries the
  type-theoretic derivation calculus.  The flattened components
  $(c, \varphi)$ and $(A, E)$ carry no intrinsic logic.

\item \textbf{Higher cohomology.}  The component-wise view records only
  $H^1$ (the obstruction class~$B$); the Cohomological Spectrum
  carries $H^n$ for all~$n$, plus the long exact sequence, spectral
  sequences, and the derived category.  These detect obstructions
  invisible to~$H^1$ alone.

\item \textbf{Residuation.}  The Graded Evaluation carries the
  residual $\multimap$, which answers ``what quality is achievable
  given the current evidence?''  The flattened grade~$T$ is a bare
  number with no such algebraic structure.

\item \textbf{Fibration structure.}  The judgment fibration
  $\pi : \JJ \to \mathbf{STop}$ (Chapter~\ref{ch:fibration}) provides
  a unified geometric framework---fibres, sections, descent,
  optimisation---that has no counterpart in the flat component-wise
  view.
\end{enumerate}

These are not additional axioms but \emph{consequences} of the
structured formulation: they are present in the ambient mathematical
structures and become available in the 4-tuple.  The equivalence
theorem guarantees backward compatibility---no theorem is lost---while
the structured view provides forward capability.
\end{remark}
